% %%%%%%%%%%%%%%%%%%%%%%%
% %%%%%%%%%%%%%%%%%%%%%%%
% % ICML format
% %%%%%%%% ICML 2025 EXAMPLE LATEX SUBMISSION FILE %%%%%%%%%%%%%%%%%

% \documentclass{article}

% % Recommended, but optional, packages for figures and better typesetting:
% \usepackage{microtype}
% \usepackage{graphicx}
% % \usepackage{subfigure}
% \usepackage{booktabs} % for professional tables

% % hyperref makes hyperlinks in the resulting PDF.
% % If your build breaks (sometimes temporarily if a hyperlink spans a page)
% % please comment out the following usepackage line and replace
% % \usepackage{icml2025} with \usepackage[nohyperref]{icml2025} above.
% \usepackage{hyperref}


% % Attempt to make hyperref and algorithmic work together better:
% \newcommand{\theHalgorithm}{\arabic{algorithm}}

% % Use the following line for the initial blind version submitted for review:
% \usepackage{icml2025}

% % If accepted, instead use the following line for the camera-ready submission:
% % \usepackage[accepted]{icml2025}


%%%%%%%%%%%%%%%%%%%%%%
%%%%%%%%%%%%%%%%%%%%%%
% ACL ARR format
% This must be in the first 5 lines to tell arXiv to use pdfLaTeX, which is strongly recommended.
\pdfoutput=1
% In particular, the hyperref package requires pdfLaTeX in order to break URLs across lines.

\documentclass[11pt]{article}

% Change "review" to "final" to generate the final (sometimes called camera-ready) version.
% Change to "preprint" to generate a non-anonymous version with page numbers.
\usepackage[review]{acl}

%%%%%%%%%%%%%%%%%%%%%%%
%%%%%%%%%%%%%%%%%%%%%%%
% % ICLR format
% \documentclass{article} % For LaTeX2e
% \usepackage{iclr2024_conference,times}

% % Optional math commands from https://github.com/goodfeli/dlbook_notation.
% \input{math_commands.tex}

% \usepackage{hyperref}
% \usepackage{url}

% Standard package includes
\usepackage{times}
\usepackage{latexsym}

% % For proper rendering and hyphenation of words containing Latin characters (including in bib files)
% \usepackage[T1]{fontenc}
% % For Vietnamese characters
% % \usepackage[T5]{fontenc}
% % See https://www.latex-project.org/help/documentation/encguide.pdf for other character sets

% % This assumes your files are encoded as UTF8
% \usepackage[utf8]{inputenc}

% % This is not strictly necessary, and may be commented out,
% % but it will improve the layout of the manuscript,
% % and will typically save some space.
% \usepackage{microtype}

% % This is also not strictly necessary, and may be commented out.
% % However, it will improve the aesthetics of text in
% % the typewriter font.
% \usepackage{inconsolata}

% %Including images in your LaTeX document requires adding
% %additional package(s)
% \usepackage{graphicx}

% If the title and author information does not fit in the area allocated, uncomment the following
%
%\setlength\titlebox{<dim>}
%
% and set <dim> to something 5cm or larger.

% self-imported packages
\usepackage{cleveref}
% \usepackage{algorithm2e}
\usepackage{utfsym}
\usepackage{fontawesome}
\usepackage{multirow}
\usepackage{booktabs}
\usepackage{adjustbox}
\usepackage{enumitem}
\usepackage{graphicx}
\usepackage{xcolor}
\usepackage{colortbl}
\usepackage{pgf}
\usepackage{pgfplotstable} % For handling and displaying tables
\usepackage{colortbl} % For coloring table cells
\usepackage{booktabs} % For better-looking tables
\usepackage{etoolbox}
\usepackage{twemojis}
\usepackage{wrapfig,lipsum,booktabs}
\usepackage{graphicx}
\usepackage{caption}
\usepackage{amssymb}% http://ctan.org/pkg/amssymb
\usepackage{pifont}% http://ctan.org/pkg/pifont
\usepackage{skak}
\usepackage{accents}
\usepackage{subcaption}

\usepackage[normalem]{ulem}
\useunder{\uline}{\ul}{}

% self-defined command
\newcommand{\dnaihao}[1]{\textcolor{blue}{[naihao: #1]}}
\newcommand{\zshe}[1]{\textcolor{red}{[sheng: #1]}}
\newcommand{\shuaichen}[1]{\textcolor{orange}{[shuaichen: #1]}}
\newcommand{\patricng}[1]{\textcolor{teal}{[PN: #1]}}
\newcommand{\yiqun}[1]{\textcolor{violet}{[yiqun: #1]}}
% comment to remove comment
%\renewcommand{\zshe}[1]{}
%\renewcommand{\dnaihao}[1]{}
%\renewcommand{\shuaichen}[1]{}
%\renewcommand{\patricng}[1]{}
\newcommand{\cmark}{\ding{51}}%
\newcommand{\xmark}{\ding{55}}%
\newcommand{\modelname}[0]{\textbf{\textit{TableMistral}}}
\newcommand{\grayc}[0]{\cellcolor[rgb]{0.957,0.957,0.957}}
\definecolor{mygray}{gray}{0.75}
\newcommand{\g}[0]{\color{mygray}}
\definecolor{myred}{rgb}{0.75, 0.0, 0.0}
\newcommand{\red}[0]{\color{myred}}
\newcommand{\crown}[0]{
\includegraphics[height=1em]{imgs/crown.pdf}}



\usepackage{listings}
\usepackage{enumitem}
\lstset{
basicstyle=\small\ttfamily,
columns=flexible,
breakindent=0pt,
breaklines=true
}
\usepackage{wrapfig,lipsum,booktabs}

% \title{Towards Understanding Effective Strategies for building TableLLM Models}
% \title{Beyond SOTA: Towards Better Understanding\\the Limitations of Existing Table Instruction Tuning}
% \title{Towards Decoupling Training Data and Base Model Effects\\in Table Instruction Tuning}
% \title{Towards Better Understanding Table Instruction Tuning:\\Decoupling the Effects from Data versus Models}

\title{Towards Better Understanding Table Instruction Tuning: Decoupling the Effects from Data versus Models}

% \twocolumn[
% \icmltitle{Towards Better Understanding Table Instruction Tuning:\\Decoupling the Effects from Data versus Models}

% Author information can be set in various styles:
% For several authors from the same institution:
% \author{Author 1 \and ... \and Author n \\
%         Address line \\ ... \\ Address line}
% if the names do not fit well on one line use
%         Author 1 \\ {\bf Author 2} \\ ... \\ {\bf Author n} \\
% For authors from different institutions:
% \author{Author 1 \\ Address line \\  ... \\ Address line
%         \And  ... \And
%         Author n \\ Address line \\ ... \\ Address line}
% To start a separate ``row'' of authors use \AND, as in
% \author{Author 1 \\ Address line \\  ... \\ Address line
%         \AND
%         Author 2 \\ Address line \\ ... \\ Address line \And
%         Author 3 \\ Address line \\ ... \\ Address line}

% \iclrfinalcopy

\author{Naihao Deng$^{1}$\thanks{Work done during internship at AWS AI Labs}, Sheng Zhang$^{2}$, Henghui Zhu$^{2}$, \textbf{Shuaichen Chang}$^{2}$, \textbf{Jiani Zhang}$^{2}$,\\
\textbf{Alexander Hanbo Li}$^{2}$, \textbf{Chung-Wei Hang}$^{2}$, \textbf{Hideo Kobayashi}$^{2}$, \textbf{Yiqun Hu}$^{2}$, \textbf{Patrick Ng}$^{2}$ \\
$^{1}$University of Michigan \ $^{2}$AWS AI Labs, New York, USA\\
$^{1}$\texttt{dnaihao@umich.edu}; 
% \texttt{\{zshe, patricng, henghui, hanboli, juwanga, yiqunhu, wyw, zhiguow, bxiang\}@amazon.com} \\
$^{2}$\texttt{\{zshe, henghui, patricng\}@amazon.com} \\
}

%\author{
%  \textbf{First Author\textsuperscript{1}},
%  \textbf{Second Author\textsuperscript{1,2}},
%  \textbf{Third T. Author\textsuperscript{1}},
%  \textbf{Fourth Author\textsuperscript{1}},
%\\
%  \textbf{Fifth Author\textsuperscript{1,2}},
%  \textbf{Sixth Author\textsuperscript{1}},
%  \textbf{Seventh Author\textsuperscript{1}},
%  \textbf{Eighth Author \textsuperscript{1,2,3,4}},
%\\
%  \textbf{Ninth Author\textsuperscript{1}},
%  \textbf{Tenth Author\textsuperscript{1}},
%  \textbf{Eleventh E. Author\textsuperscript{1,2,3,4,5}},
%  \textbf{Twelfth Author\textsuperscript{1}},
%\\
%  \textbf{Thirteenth Author\textsuperscript{3}},
%  \textbf{Fourteenth F. Author\textsuperscript{2,4}},
%  \textbf{Fifteenth Author\textsuperscript{1}},
%  \textbf{Sixteenth Author\textsuperscript{1}},
%\\
%  \textbf{Seventeenth S. Author\textsuperscript{4,5}},
%  \textbf{Eighteenth Author\textsuperscript{3,4}},
%  \textbf{Nineteenth N. Author\textsuperscript{2,5}},
%  \textbf{Twentieth Author\textsuperscript{1}}
%\\
%\\
%  \textsuperscript{1}Affiliation 1,
%  \textsuperscript{2}Affiliation 2,
%  \textsuperscript{3}Affiliation 3,
%  \textsuperscript{4}Affiliation 4,
%  \textsuperscript{5}Affiliation 5
%\\
%  \small{
%    \textbf{Correspondence:} \href{mailto:email@domain}{email@domain}
%  }
%}

\begin{document}
\maketitle

\begin{abstract}

Recent advances in natural language processing have leveraged instruction tuning to enhance Large Language Models (LLMs) for table-related tasks. 
However, previous works train different base models with different training data, lacking an apples-to-apples comparison across the result table LLMs. 
To address this, we fine-tune base models from the Mistral, OLMo, and Phi families on existing public training datasets. 
% We find that while these models achieve competitive or even state-of-the-art (SOTA) performance on the data that they are trained on, they suffer issues of poor out-of-domain performance and the loss of general capabilities compared to the base model.
% \dnaihao{not a single model can consistently outperform their base model.}
% \patricng{Why do we want to highlight this is only for Mistral model in the abstract? If our story is on model bias, we should at least have 3 base FMs. (I do see OLMo in Table2)}
Our replication achieves performance on par with or surpassing existing table LLMs, establishing new state-of-the-art performance on Hitab, a table question-answering dataset.
More importantly, through systematic out-of-domain evaluation, we decouple the contributions of training data and the base model, providing insight into their individual impacts.
In addition, we assess the effects of table-specific instruction tuning on general-purpose benchmarks, revealing trade-offs between specialization and generalization.

\end{abstract}

\section{Introduction}
Researchers have tried to improve table understanding abilities of Large Language Models (LLMs) through instruction tuning on table tasks, aiming to construct a generalist model for table understanding.
Existing work adapts table understanding benchmarks for table instruction tuning \citep{zhang-etal-2024-tablellama, zheng2024multimodal} or synthesize instruction-answer pairs using LLMs to fine-tune smaller scale models \citep{li2023table, zhang2024tablellm, wu2024tablebenchcomprehensivecomplexbenchmark}.

% However, progress in this area is difficult to measure due to the lack of apples-to-apples comparisons across studies.
However, these studies utilize different pretrained model architectures, training datasets and evaluation benchmarks, making it difficult to determine the specific sources of their contributions.
% However, different works use varied foundational models and training datasets, making it challenging to isolate the contributions of the base models from those of the training data.
For example, TableLlama \citep{zhang-etal-2024-tablellama} is trained from LongLoRA \cite{longlora}, a variant of Llama 2 \citep{touvron2023llama}. 
TableLLM \cite{zhang2024tablellm} is based on CodeLlama Instruct. 
TableBenchLLM \cite{wu2024tablebenchcomprehensivecomplexbenchmark} employs models such as Llama 3.1 \cite{dubey2024llama} and Qwen2 \cite{yang2024qwen2}, and models such as Table-GPT \citep{li2023table} are closed-source trained on GPT-3.5 \cite{ouyang2022training}.
In addition, these models use different training datasets, with some relying on existing benchmarks (e.g., TableLlama) and others introducing proprietary training datasets (e.g., Table-GPT).
Furthermore, previous work chooses different sets of evaluation benchmarks and seldomly compares their models with others, leaving an unclear image of how much progress researchers have made in building table LLMs.

For a fair and comprehensive study, we fine-tune the same base models from the Mistral \citep{jiang2023mistral}, OLMo \citep{groeneveld2024olmo}, and Phi \citep{abdin2024phi} families using the respective training data from each work.
Our experiments demonstrate that these fine-tuned models perform comparably to or outperform current table LLMs on their respective evaluation datasets. 
Notably, our models establish new state-of-the-art (SOTA) results on HiTab, a table question answering benchmark.
Beyond replication, we conduct systematic evaluations on diverse table tasks, including table question answering (Table QA), table-to-text generation, table fact verification, and beyond.
These tasks span real-world datasets \citep{nan-etal-2022-fetaqa, cheng-etal-2022-hitab} and synthetic data proposed in existing works \citep{li2023table, wu2024tablebenchcomprehensivecomplexbenchmark}.
To assess trade-offs, we also evaluate these models on general instruction following and reasoning benchmarks, analyzing how specializing in table tasks affects their general capabilities.
Our work disentangles the effects of training data and base models.

Specifically, we find that 1) The base models demonstrate competitive performance on table benchmarks compared to the fine-tuned models; 
2) Certain training data (e.g. the training data from TableLLM) consistently outperforms the others on the out-of-domain table tasks;
3) SOTA-chasing can be meaningless as the model's performance does not generalize to datasets in the same task category;
4) Transferability exists across table tasks;
5) The effects of the training data depend on the choice of base models, and strong base models lead to better performance;
6) Proper fine-tuning does not necessarily compromise the model's general capabilities.
% We reveal that one set of training data may have different performance with a different base model, 
We hope our findings provide actionable insights into model selection and dataset construction for building effective table LLMs.

\noindent In summary, our contributions are three fold.

\begin{itemize}[leftmargin=0.4cm,itemsep=0.1cm]
    \item We replicate existing table LLMs by fine-tuning various base models, achieving comparable or superior performance on benchmarks reported in the existing works, respectively.
    % and setting new SOTA results on HiTab.
    \item We decouple the contributions of the training data and base models, revealing that different base models perform differently with the same set of training data.
    % In addition, we offer practical insights into data construction, model selection, and strategies for building competitive table LLMs.
    We provide our findings on the effects of training data and base models.
    \item We expand the evaluation topology of table LLMs to general benchmarks, revealing that proper fine-tuning does not necessarily compromise the model's general capabilities.
    % Moreover, we provide the first open-source implementation of Table-GPT, one of the existing closed-source table LLMs. \patricng{I think this bullet point can be split into at least 2 different contribution items}
    % \patricng{Add examples of what you mean by evaluation topology here}
    % \item Our findings reveal the limitations of the current table instruction tuning, that none of the models trained on the existing table instruction tuning data can consistently outperform their base model on out-of-domain table tasks, and that they decay in general capabilities such as instruction following ability, inviting future efforts in our field to construct an authentic generalist table LLM. 
\end{itemize}

\section{Related Works}

\paragraph{Table-Related Tasks.}


Earlier work has focused on extracting table content from HTML \citep{chen-etal-2000-mining, tengli-etal-2004-learning}.
The deep learning era has seen more diverse table-related tasks such as table question answering (table QA), the task of answering a question given the table and certain context in the format of multiple-choice \citep{jauhar2016tabmcq} and free-form answer \citep{nan-etal-2022-fetaqa}; 
table fact verification, the task of determining whether a given claim is supported or refuted by the table content \citep{chen2019tabfact, gupta-etal-2020-infotabs}; 
table-to-text, the task of generating a description given the table or some highlighted table cells \citep{parikh-etal-2020-totto}.
These proposed benchmarks cover a diverse set of domains, including Wikipedia tables \citep{parikh-etal-2020-totto}, financial tables \citep{chen-etal-2021-finqa}, scientific tables \citep{moosavi2021scigen}, which serve as invaluable sources for developing and testing general table understanding models.

% Moreover, there are many other tasks that are relevant to table understanding.
% For instance, \citet{xie-etal-2022-unifiedskg} have demonstrated the synergy effects of models learning on table understanding tasks and other structured knowledge tasks tasks such as SQL generation \citep{yu-etal-2018-spider}, program to text \citep{shu-etal-2021-logic}.
% Specifically, tasks such as SQL generation also requires table understanding abilities to link relevant entities from database table and the natural language query \citep{deng-etal-2022-recent}.



\paragraph{Methods for Table Understanding.}

Researchers have explored various methods for table understanding in the past decade such as adapting the model's internal structures to align with table structures \citep{lebret-etal-2016-neural, liu2018table, yang-etal-2022-tableformer}, synthesizing a large table pre-training corpus and designing table-specific training objectives \citep{yin-etal-2020-tabert, herzig-etal-2020-tapas}.
% With the advent of the sequence-to-sequence (seq2seq) framework \citep{sutskever2014sequence}, researchers try to impart structured knowledge to the seq2seq models by adapting their internal structures \citep{lebret-etal-2016-neural, liu2018table}.
% This remains to be one of the directions in the transformer era \citep{vaswani2017attention} where researchers try to adapt the attention mechanism to align with the table structure \citep{yang-etal-2022-tableformer}.
% With the success of pre-training methods \citep{devlin-etal-2019-bert}, researchers have also started a data-driven methods to improve models' table understanding ability by synthesizing large table pre-training corpus and designing table-specific training objectives \citep{yin-etal-2020-tabert, herzig-etal-2020-tapas}.
Recently, the LLM era has witnessed a paradigm shift for research on tables.
As LLMs have inherent abilities on table understanding, researchers employ prompt engineering on these LLMs for better performance on tables \citep{chang2023prompt, deng-etal-2024-tables}.
Another line of research involves instruction tuning LLMs by adapting existing table-related benchmarks.
This leads to various table LLMs such as Table-GPT \citep{li2023table}, TableLlama \citep{zhang-etal-2024-tablellama}, TableLlava \citep{zheng2024multimodal}, and TableLLM \citep{zhang2024tablellm}. 
Among them, while TableLlama achieves decent performance on in-domain data, it suffers a significant performance drop on unseen table tasks \citep{zheng2024multimodal}.
\citet{su2024tablegpt2} introduce TableGPT-2, a concurrent work to ours that utilizes synthesized training data for model training. 
Due to the overlapping timelines, we do not include their model in our study.
% Moreover, TableLlama and TableLlava have consumed significant amount of data in the instruction tuning stage.
Though representing tables as images is a promising direction, recent works reveal that its performance still lags behind representing tables as texts \citep{deng-etal-2024-tables}.
Therefore, in this paper, we focus on the text representation of tables.

\begin{table}[t]
    \small
    \centering
    \renewcommand{\arraystretch}{1.3}
    \setlength\tabcolsep{5pt}
    % \scalebox{1}{
    \resizebox{\linewidth}{!}{
    \begin{tabular}{cccccccc}
        \toprule
        Model & Base Model & \begin{tabular}[c]{@{}c@{}}Data\\Size\end{tabular} & \begin{tabular}[c]{@{}c@{}}Data\\Source\end{tabular} & \begin{tabular}[c]{@{}c@{}}Open\\Model?\end{tabular} & \begin{tabular}[c]{@{}c@{}}Open\\Data?\end{tabular}\\
        \midrule    
        TableGPT (\citeyear{zha2023tablegpt}) & - & - & - & \xmark & \xmark \\
        \grayc Table-GPT (\citeyear{li2023table}) & \grayc GPT-3.5 & \grayc 66K & \grayc S & \grayc \xmark & \grayc \cmark\\
        TableLlama (\citeyear{zhang-etal-2024-tablellama}) &  LongLoRA 7B$^\dagger$ & 2M & R & \cmark & \cmark\\
        \grayc TableLLM (\citeyear{zhang2024tablellm}) & \grayc \begin{tabular}[c]{@{}c@{}}CodeLlama 7B\\\& 13B Instruct\end{tabular} & \grayc 80.5K & \grayc R + S & \grayc \cmark & \grayc \cmark\\
        \begin{tabular}[c]{@{}c@{}}TableBenchLLM\\ (\citeyear{wu2024tablebenchcomprehensivecomplexbenchmark})\end{tabular} & \begin{tabular}[c]{@{}c@{}}Llama 3.1-8B\\\& others\end{tabular} & 20K & S & \cmark & \cmark\\
         \bottomrule
    \end{tabular}
    }
    \caption{Information for current table instruction tuned models.
    For the ``Data Source'', ``S'' represents synthesized data, and ``R'' represents the real data.
    $\dagger$: TableLlama has adopted the LongLoRA \cite{longlora} base model, which is a variant based on the Llama 2 7B model with a longer context window.
    }
    \label{tab: tablellm-hyperparameter}
\end{table}

% \section{Training Table LLMs on the Same Base Model}
\section{Replicating Existing Table LLMs}


\begin{table*}[t]
    \centering
    \small
    \renewcommand{\arraystretch}{1.3}
    \setlength\tabcolsep{2pt}
    \scalebox{1}{
    % \resizebox{\linewidth}{!}{
    \begin{tabular}{c|ccc|ccccccc}
    \toprule
     \multicolumn{1}{c|}{\multirow{2}{*}{Base Models}} & FeTaQA  &  HiTab & TabFact & FEVEROUS  &  HybridQA & KVRET & ToTTo & WikiSQL & WikiTQA \\
    & (BLEU) & (Acc) & (Acc) & (Acc) & (Acc) & (F1\textsubscript{Micro}) & (BLEU) & (Acc) & (Acc)\\
    \midrule
    % \begin{tabular}[c]{@{}c@{}}TableLlama (\citeyear{zhang-etal-2024-tablellama})\end{tabular} & \textbf{39.05} & 64.71 & 82.55 & 73.77 & \textbf{39.38} & \underline{48.73} & 20.77 & 50.48 & 35.01 \\
    % \midrule
    % \begin{tabular}[c]{@{}c@{}}Mistral-TableLlama \end{tabular}  & \underline{38.66} & \textbf{70.58}$^\dagger$ & \textbf{86.83} & \underline{75.94} & 27.18 & 46.65 & \underline{28.54} & \textbf{64.50} & \underline{47.38}\\
    % % Mistral\textsubscript{Nemo}-TableLlama & 39.37 & 64.46 & 89.51 & 81.52 & 32.49 & 45.52 & 33.61 & 66.54 & 44.64 \\
    % OLMo-TableLlama & 36.79 & \underline{67.93} & 83.75 & 69.78 & 20.31 & 44.62 & 20.82 & 56.90 & 38.81 \\
    % Phi-TableLlama & 38.06 & 63.57 & \underline{86.19} & \textbf{78.30} & \underline{33.64} & \textbf{56.03} & \textbf{29.57} & \underline{63.30} & \textbf{47.70} \\
  \multicolumn{10}{l}{\grayc \textit{Original \citep{zhang-etal-2024-tablellama}}}\\
  LongLoRA 7B$^\ddagger$ & \textbf{39.0} & 64.7 & 82.5 & 73.8 & \textbf{39.4} & \underline{48.7} & 20.8 & 50.5 & 35.0 \\
\midrule
\multicolumn{10}{l}{\grayc \textit{Ours}}\\
\multicolumn{1}{c|}{Mistral v0.3 7B Instruct}  & \underline{38.7} & \textbf{70.6}$^\dagger$ & \textbf{86.8} & \underline{75.9} & 27.2 & 46.6 & \underline{28.5} & \textbf{64.5} & \underline{47.4}\\
% Mistral\textsubscript{Nemo}-TableLlama & 39.4 & 64.5 & 89.5 & 81.5 & 32.5 & 45.5 & 33.6 & 66.5 & 44.6 \\
\multicolumn{1}{c|}{OLMo 7B Instruct} & 36.8 & \underline{67.9} & 83.8 & 69.8 & 20.3 & 44.6 & 20.8 & 56.9 & 38.8 \\
\multicolumn{1}{c|}{Phi 3 Small Instruct (7B)} & 38.1 & 63.6 & \underline{86.2} & \textbf{78.3} & \underline{33.6} & \textbf{56.0} & \textbf{29.6} & \underline{63.3} & \textbf{47.7} \\

    \bottomrule
    \end{tabular}
    }
    \caption{Performance comparison between the original TableLlama and our fine-tuned models from different model families on the in-domain tuned (left three columns) and out-of-domain (right six columns) datasets.
    The number is bold if it is the best among the four, and underscored if it is the second.
    $\dagger$: we surpass the previous SOTA performance (64.7 by TableLlama) on HiTab.
    }
    \label{tab:mistral-tablellama}
\end{table*}

\paragraph{Issues with Comparing Existing Table LLMs}


% There are two issues with the current table LLMs.


% \begin{table}[t]
%     \centering
%     \small
%     \renewcommand{\arraystretch}{1.3}
%     \setlength\tabcolsep{2pt}
%     \resizebox{\linewidth}{!}{
%     \begin{tabular}{p{2.3cm}ccc|cc}
%     \toprule
%     Model & Company & \begin{tabular}[c]{@{}l@{}}Open\\Model?\end{tabular} & \begin{tabular}[c]{@{}l@{}}Open\\License?\end{tabular} &  MMLU & HellaSwag\footnotemark[1]\\
%     \midrule
%     Llama 3.1 Instruct 8B (\citeyear{dubey2024llama})  &  \multirow{2}{*}{Meta} & \multirow{2}{*}{\cmark} & \multirow{2}{*}{\xmark\footnotemark[2]} & \multirow{2}{*}{66.04} & \multirow{2}{*}{80.01} \\
%     Mistral v0.3 7B Instruct\footnotemark[3] & \multirow{2}{*}{Mistral AI} & \multirow{2}{*}{\cmark} & \multirow{2}{*}{\cmark} & \multirow{2}{*}{61.37} & \multirow{2}{*}{84.84}\\
%     \bottomrule
%     \end{tabular}
%     }
%     \caption{Comparisons of the Mistral v0.3 7B Instruct model versus the Llama 3.1 Instruct 8B model in terms of their license and their performance on the general MMLU benchmark (5-shot) and HellaSwag (10-shot), a grade school math dataset.}
%     \label{tab:model-comparison}
% \end{table}



% \paragraph{Llama Dominance.}
% In \Cref{tab: tablellm-hyperparameter}, we note that all of the open-source models trained by the community are based from Llama models.
% However, as depicted in \Cref{tab:model-comparison}, though Llama models are open models, they do not have an open license\footnotemark[2].
% It is restricted to research and non-commercial use only, therefore ruling out the contribution from the industry research commuinty.
% In addition, the license restricts the re-distribution of the Llama models.
% If researchers modify or improve the model, they cannot legally share those modifications with others unless it adheres to the limited terms of the license.

% To avoid the potential license issue, we employ Mistral v0.3 7B Instruct model, the most advanced model from Mistral family with an open license.
% In addition, as shown in \Cref{tab:model-comparison}, Mistral v0.3 7B Instruct model yields comparable results with Llama 3.1 Instruct 8B on the general MMLU benchmark \cite{hendrycks2020measuring} and HellaSwag \cite{zellers-etal-2019-hellaswag}, a grade school math dataset.

\Cref{tab: tablellm-hyperparameter} outlines the base models used in existing table LLMs. 
% \shuaichen{This paragraph is too short. can we merge it with the next paragraph}
% which are fine-tuned from \textit{different foundational models using different datasets.} 
% \shuaichen{trained -$>$ finetuned/adapted?}
These base models, ranging from various Llama models to closed-source models such as GPT-3.5, differ significantly in their architecture designs, model sizes, and training recipes.
In addition, each table LLM introduces its own unique training data, making it challenging to disentangle the impact of the training data from that of the base model.

To explore the contribution of the training data used in existing table LLM works, we train the same base models on datasets utilized in each of the existing works.
We first demonstrate that our implementation yields comparable or better results than the performance reported in the existing works in \Cref{sec: reimpl-compare}. We then evaluate our trained models across various setups in \Cref{sec: real-world-eval,sec: general-eval}.

% \subsection{Experimental Setup}

\paragraph{Foundational LLM Selections.}
For the training data from each existing work, we fine-tune \texttt{Mistral-7B-Instruct-v0.3} \citep{jiang2023mistral}, OLMo 7B Instruct \citep{groeneveld2024olmo} and Phi 3 Small Instruct (7B) \citep{abdin2024phi}.
Following \citet{zhang-etal-2024-tablellama, zhang2024tablellm, wu2024tablebenchcomprehensivecomplexbenchmark}, we fine-tune all the models through full parameter fine-tuning.

\paragraph{Experimental Setups.}
To rule out the effects of the learning rate, we train all three models using a set of learning rates: 5e-5, 1e-5, 5e-6, 1e-6, 5e-7, 1e-7, 5e-8, and 1e-8.
Empirically, we find that they achieve the best when the learning rate is 5e-7.
%Therefore, we train all of our models at the learning rate of 5e-7.
% \paragraph{Number of Epochs.}
We do not see significant performance changes as we increase the training steps.
For consistency, we fine-tune our models for three epochs across all the experiments.
More details in \Cref{app-sec: experimental-setup}.

\section{Comparison of Our Models v.s. the Original Models}
\label{sec: reimpl-compare}

Here we report the performance of our fine-tuned models based on Mistral v0.3 7B Instruct, OLMo 7B Instruct, and Phi 3 Small Instruct (7B) versus the original models on the datasets reported in each of the original work.

\subsection{Replicating TableLlama}
\label{subsec: reimpl-tablellama}

\paragraph{Training Datasets.} 
The original TableLlama \citep{zhang-etal-2024-tablellama} uses 2 million data points in its instruction tuning stage, which can be unnecessarily large.
In addition, we do not have enough computing resources to instruction-tune our model on a dataset of such a scale.
Therefore, we rule out the table operation datasets and only maintain the training data for FeTaQA \citep{nan-etal-2022-fetaqa}, HiTab \citep{cheng-etal-2022-hitab}, and TabFact \citep{chen2019tabfact} to fine-tune our model, which results in 107K training instances.

\paragraph{Evaluation Datasets.}
Following \citet{zhang-etal-2024-tablellama}, we use the FeTaQA \citep{nan-etal-2022-fetaqa}, HiTab \citep{cheng-etal-2022-hitab}, and TabFact \citep{chen2019tabfact} as the in-domain evaluation sets.
In addition, we compare our fine-tuned models versus the original TableLlama on FEVEROUS \citep{aly-etal-2021-fact}, HybridQA \citep{chen-etal-2020-hybridqa}, KVRET \citep{eric2017key}, ToTTo \citep{parikh-etal-2020-totto}, WikiSQL \citep{zhong2017seq2sql}, and WikiTQ \citep{pasupat-liang-2015-compositional}.

\paragraph{Comparison.}
\Cref{tab:mistral-tablellama} compares the original TableLlama model (first row) versus our fine-tuned models.
Our fine-tuned models yield similar or better performance than the original TableLlama model in most cases.
In addition, we achieve the new SOTA performance on HiTab by fine-tuning the Mistral model.
As we only use 107K (5\% of the 2M data points used by the original TableLlama), our results demonstrate that \textit{with proper instruction-tuning, we can achieve competitive results on table tasks with much fewer data.}
% \zshe{Do we need to add the current SOTA in Table 2? It may be confusing that LongLoRA 7B is the SOTA. }

\subsection{Replicating TableLLM}

\begin{table}[t]
    \centering
    \small
    \renewcommand{\arraystretch}{1.3}
    \setlength\tabcolsep{2pt}
    \scalebox{1}{
    % \resizebox{\linewidth}{!}{
    \begin{tabular}{c|ccc|c}
    \toprule
     Base &  WikiTQ\textsubscript{m} & TATQA\textsubscript{m} & FeTaQA\textsubscript{m} & OTT-QA\textsubscript{m} \\
    Models & (Acc\textsubscript{p}) & (Acc\textsubscript{p}) & (BLEU) & (Acc\textsubscript{p}) \\
    \midrule
    % \begin{tabular}[c]{@{}c@{}}TableLLM (\citeyear{zhang2024tablellm})\end{tabular} & 72.51 & 51.13 & 8.41 & 57.27 \\
    % \midrule
    % \begin{tabular}[c]{@{}c@{}}Mistral-TableLLM\end{tabular} & \textbf{75.99} & \underline{55.38} & \underline{10.60} & \textbf{64.32}\\
    % % Mistral\textsubscript{Nemo}-TableLLM & 79.94 & 58.13 & 11.60 & 67.39 \\
    % OLMo-TableLLM & 66.82 & 50.25 & 10.53 & 58.08 \\
    % Phi-TableLLM & \underline{75.36} & \textbf{57.75} & \textbf{12.10} & \underline{63.31} \\
      \multicolumn{5}{l}{\grayc \textit{Original \citep{zhang2024tablellm}}}\\
    CodeLlama$^\ddagger$ & 72.5 & 51.1 & 8.4 & 57.3 \\
\midrule
\multicolumn{5}{l}{\grayc \textit{Ours}}\\
\multicolumn{1}{c|}{Mistral} & \textbf{76.0} & \underline{55.4} & \underline{10.6} & \textbf{64.3}\\
% Mistral\textsubscript{Nemo}-TableLLM & 79.9 & 58.1 & 11.6 & 67.4 \\
\multicolumn{1}{c|}{OLMo} & 66.8 & 50.2 & 10.5 & 58.1 \\
\multicolumn{1}{c|}{Phi} & \underline{75.4} & \textbf{57.8} & \textbf{12.1} & \underline{63.3} \\
    \bottomrule
    \end{tabular}
    }
    \caption{Performance comparison between the original TableLLM and our fine-tuned models.
    % on the in-domain tuned (left three columns) and out-of-domain (rightmost column) datasets.
    All four models are 7B and instruction-tuned.
    We denote the evaluation datasets with a subscript ``m'' as they are adapted by \citet{zhang2024tablellm}.
    }
    \label{tab:mistral-tablellm}
\end{table}


\paragraph{Training Datasets.} 
We use the original instruction-tuning set by \citet{zhang2024tablellm}, which includes 80.5K training instances.

\paragraph{Evaluation Datasets.}
Following \citet{zhang2024tablellm}, we use the modified version of WikiTQ \citep{pasupat-liang-2015-compositional}, TATQA \citep{zhu-etal-2021-tat}, and FeTaQA \citep{nan-etal-2022-fetaqa} as the in-domain evaluation sets, and OTT-QA \citep{chen2020open} as the out-of-domain evaluation set.

\paragraph{Comparison.}
\Cref{tab:mistral-tablellm} compares the original TableLLM versus our fine-tuned models.
We note that our evaluation metrics are distinct from what \citet{zhang2024tablellm} have used originally.
\citet{zhang2024tablellm} use CritiqueLLM \citep{ke-etal-2024-critiquellm} as a judge to decide the correctness of the answers.
However, the model judgments are made in Chinese\footnote{\citet{zhang2024tablellm}'s inference results are available at \url{https://github.com/RUCKBReasoning/TableLLM/blob/main/inference/results/TableLLM-7b/Grade\_fetaqa.jsonl}}, a different language from the language in all the training and evaluation datasets.
In addition, the scores assigned by the CritiqueLLM is not consistent for a single evaluation example.
Therefore, for WikiTQ\textsubscript{m}, TATQA\textsubscript{m}, and OTT-QA\textsubscript{m}, we report the Acc\textsubscript{p} scores, where we calculate whether the gold answer entities appear in the model's response.
We find that our fine-tuned models based on the Mistral and Phi models consistently outperform the original TableLLM model on these datasets, and we attribute the performance improvement to the stronger base model (Mistral v0.3 7B Instruct and Phi 3 Small Instruct) we have versus theirs (CodeLlama 7B Instruct).


\subsection{Replicating TableBenchLLM}

\begin{table}[t]
    \centering
    \small
    \renewcommand{\arraystretch}{1.3}
    \setlength\tabcolsep{2pt}
    \scalebox{1}{
    % \resizebox{\linewidth}{!}{
    \begin{tabular}{c|cc}
    \toprule
    % \midrule
    % Fact Checking & 63.68 & 49.38\\
    % Num-Reasoning & 18.01 & 14.78\\
    % Data Anlysis & 28.32 & 30.82\\
    % Visualization & 24.00 & 1.21\\
    Base & TableBench\textsubscript{eval} \\
    Models & (R-L) \\
     \midrule
    % \begin{tabular}[c]{@{}c@{}}TableBenchLLM\textsubscript{Llama-3.1-8B} (\citeyear{wu2024tablebenchcomprehensivecomplexbenchmark})\end{tabular} & 27.19 \\
    \multicolumn{2}{l}{\grayc \textit{Original \citep{wu2024tablebenchcomprehensivecomplexbenchmark}}}\\
   Llama 3.1 8B $^\ddagger$ & \underline{27.2} \\ 
    \midrule
    % \begin{tabular}[c]{@{}c@{}}Mistral-TableBenchLLM \end{tabular} & \underline{27.22}\\    
    \multicolumn{2}{l}{\grayc \textit{Ours}}\\
    \multicolumn{1}{c|}{Mistral v0.3 7B Instruct} & \underline{27.2} \\
    % Mistral\textsubscript{Nemo}-TableBenchLLM & 30.51 \\
    % OLMo-TableBenchLLM & 19.31 \\
    \multicolumn{1}{c|}{OLMo 7B Instruct} & 19.3 \\
    % Phi-TableBenchLLM & \textbf{27.76} \\
    \multicolumn{1}{c|}{Phi 3 Small Instruct (7B)} & \textbf{27.8} \\
    \bottomrule
    \end{tabular}
    }
    \caption{Performance comparison between the original TablebBenchLLM based on Llama 3.1 8B and our fine-tuned models.
    ``R-L'' denotes the ROUGE-L score.
    }
    \label{tab:mistral-tablebenchllm}
\end{table}


% \begin{table*}[t]
%     \centering
%     \small
%     \renewcommand{\arraystretch}{1.3}
%     \setlength\tabcolsep{2pt}
%     \scalebox{1}{
%     % \resizebox{\linewidth}{!}{
%     \begin{tabular}{c|cccc|cccc}
%     \toprule
%    & Beer & DeepM & SS-DI  & SS-Real & C & SS-CF & WikiTQ & Efth\\
%    & (F1) & (Recall) & (Acc) & (F1) & (F1) & (Acc) & (Acc) & (F1)\\
%     \midrule
%     \begin{tabular}[c]{@{}c@{}}Table-GPT (\citeyear{li2023table})\end{tabular} & 72.70 & \textbf{100.00} & \textbf{55.80} & \textbf{56.50} & 29.40 & \textbf{71.30} & \textbf{48.60} & \textbf{88.60} \\ 
%     % GPT-3.5$^\dagger$ & 50.00 & \textbf{100.00} & 42.30 & 5.80 & \underline{26.10} & \underline{46.10} & \underline{45.00} & \underline{75.70}\\
%     \midrule
%     \begin{tabular}[c]{@{}c@{}}Mistral-TableGPT\end{tabular} & \textbf{100.00} & 97.96 & 46.35 & 46.02 & 23.75 & 25.33 & 25.46 & \underline{68.27} \\
%     % Mistral\textsubscript{Nemo}-TableGPT & 85.71 & 42.86 & 39.55 & 53.80 & 35.45 & 34.01 & 40.02 & 66.33 \\
%     OLMo-TableGPT & 96.15 & \textbf{100.00} & 45.40 & 35.34 & 19.95 & 29.31 & 16.41 & 62.46\\
%     Phi-TableGPT & \underline{98.90} & 98.81 & \underline{49.35} & \underline{55.37} & \underline{24.80} & \underline{45.18} & \underline{29.97} & 68.35 \\
%     \bottomrule
%     \end{tabular}
%     }
%     \caption{Performance comparison between the original Table-GPT based on GPT-3.5 (reported by \citet{li2023table}) and our fine-tuned models.
%     }
%     \label{tab:mistral-tablegpt}
% \end{table*}


\paragraph{Training Datasets.}
We use the original instruction-tuning set by \citet{wu2024tablebenchcomprehensivecomplexbenchmark}, which includes 20K training instances.

\paragraph{Evaluation Datasets.}
Following \citet{wu2024tablebenchcomprehensivecomplexbenchmark}, we only evaluate the model on their constructed test set, which we denote as TableBench\textsubscript{eval} in \Cref{tab:mistral-tablebenchllm}.

\paragraph{Comparison.}
Following \citet{wu2024tablebenchcomprehensivecomplexbenchmark}, we report the ROUGE-L score of our Mistral-TableBenchLLM.
In \Cref{tab:mistral-tablebenchllm}, we compare our model with the scores reported by \citet{wu2024tablebenchcomprehensivecomplexbenchmark} in the original paper, corresponding to the version of TableBenchLLM fine-tuned based on Llama 3.1 8B model.
Our Mistral-TableBenchLLM and Phi-TableBenchLLM achieve similar performance scores of 27.2 and 27.8, respectively, compared to the original TableBenchLLM's 27.2.

  \begin{table}[t]
    \raggedright
\small
    \renewcommand{\arraystretch}{1.3}
    \setlength\tabcolsep{2pt}
    % \scalebox{0.9}{
    \resizebox{\linewidth}{!}{
    \begin{tabular}{c|cccc|cccc}
    \toprule
   Base & Beer & DeepM & DI  & ED & C & CF & Wiki & CTA\\
   Models & (F1) & (Recall) & (Acc) & (F1) & (F1) & (Acc) & (Acc) & (F1)\\
    \midrule
    % \begin{tabular}[c]{@{}c@{}}Table-GPT (\citeyear{li2023table})\end{tabular} & 72.70 & \textbf{100.00} & \textbf{55.80} & \textbf{56.50} & 29.40 & \textbf{71.30} & \textbf{48.60} & \textbf{88.60} \\ 
    % % GPT-3.5$^\dagger$ & 50.00 & \textbf{100.00} & 42.30 & 5.80 & \underline{26.10} & \underline{46.10} & \underline{45.00} & \underline{75.70}\\
    % \midrule
    % \begin{tabular}[c]{@{}c@{}}Mistral-TableGPT\end{tabular} & \textbf{100.00} & 97.96 & 46.35 & 46.02 & 23.75 & 25.33 & 25.46 & \underline{68.27} \\
    % % Mistral\textsubscript{Nemo}-TableGPT & 85.71 & 42.86 & 39.55 & 53.80 & 35.45 & 34.01 & 40.02 & 66.33 \\
    % OLMo-TableGPT & 96.15 & \textbf{100.00} & 45.40 & 35.34 & 19.95 & 29.31 & 16.41 & 62.46\\
    % Phi-TableGPT & \underline{98.90} & 98.81 & \underline{49.35} & \underline{55.37} & \underline{24.80} & \underline{45.18} & \underline{29.97} & 68.35 \\
    \multicolumn{9}{l}{\grayc \textit{Original \citep{li2023table}}} \\
    GPT-3.5$^\ddagger$ & 72.7 & \textbf{100.0} & \textbf{55.8} & \textbf{56.5} & 29.4 & \textbf{71.3} & \textbf{48.6} & \textbf{88.6} \\ 
% GPT-3.5$^\dagger$ & 50.0 & \textbf{100.0} & 42.3 & 5.8 & \underline{26.1} & \underline{46.1} & \underline{45.0} & \underline{75.7}\\
\midrule
\multicolumn{9}{l}{\grayc \textit{Ours}} \\
\multicolumn{1}{c|}{Mistral} & \textbf{100.0} & 98.0 & 46.4 & 46.0 & 23.8 & 25.3 & 25.5 & \underline{68.3} \\
% Mistral\textsubscript{Nemo}-TableGPT & 85.7 & 42.9 & 39.5 & 53.8 & 35.5 & 34.0 & 40.0 & 66.3 \\
\multicolumn{1}{c|}{OLMo} & 96.2 & \textbf{100.0} & 45.4 & 35.3 & 19.9 & 29.3 & 16.4 & 62.5\\
\multicolumn{1}{c|}{Phi} & \underline{98.9} & 98.8 & \underline{49.4} & \underline{55.4} & \underline{24.8} & \underline{45.2} & \underline{30.0} & 68.3 \\
    \bottomrule
    \end{tabular}
    }
    \caption{Performance comparison between the original Table-GPT and our fine-tuned models.
    }
    \label{tab:mistral-tablegpt}
    \vspace{-0.5em}
  \end{table}%
  
  \begin{table}[t]
\centering
    \small
    \renewcommand{\arraystretch}{1.3}
    \setlength\tabcolsep{2pt}
    % \scalebox{1}{
    \resizebox{\linewidth}{!}{
    \begin{tabular}{c|cccc|cccc}
    \toprule
   & Beer & DeepM & DI  & ED & C & CF & Wiki & CTA \\
   & (F1) & (Recall) & (Acc) & (F1) & (F1) & (Acc) & (Acc) & (F1)\\
   \midrule
    % 13K & 98.90 & 92.86 & 45.85 & 43.79 & \textbf{29.35} & 21.17 & 29.18 & 66.75\\
    % 66K & \textbf{100.00} & \textbf{97.96} & \textbf{46.35} & \textbf{46.02} & 23.75 & \textbf{25.33} & \textbf{29.77} & \textbf{68.27} \\
    13K & 98.9 & 92.9 & 45.9 & 43.8 & \textbf{29.4} & 21.2 & 29.2 & 66.8\\
66K & \textbf{100.0} & \textbf{98.0} & \textbf{46.4} & \textbf{46.0} & 23.8 & \textbf{25.3} & \textbf{29.8} & \textbf{68.3} \\
    \bottomrule
    \end{tabular}
    }
    \caption{Performance comparison between training Mistral v0.3 7B Instruct on 13K instances versus 66K instances provided by \citet{li2023table}.
    }
    \label{tab:mistral-tablegpt-side-finding-all}
  \end{table}


\subsection{Replicating Table-GPT}
\label{subsec: mistral-tablegpt}


\paragraph{Training Dataset.}
We use the instruction-tuning dataset provided by \citet{li2023table} that contains 66K instances.

\paragraph{Evaluation Datasets.}
We select four in-domain test sets by \citet{li2023table}, Beer for entity matching, DeepM for schema matching, Spreadsheet-DI (DI) for data imputation, and Spreadsheet-Real (ED) for error detection.
Furthermore, we report the out-of-domain performance on Column-No-Separator (C) for missing value identification, Spreadsheet-CF (CF) for column finding, WikiTQ (Wiki) for table question answering, and Efthymiou (CTA) for column type annotation.

\paragraph{Comparison.}

\Cref{tab:mistral-tablegpt} reports the results.
We note that though the size of our fine-tuned models are all 7B, they achieve better performance than Table-GPT which is based on GPT-3.5 on Beer, and comparable performance on DeepM.
% On SS-DI and SS-Real, our Mistral-TableGPT outperforms the untuned GPT-3.5 model.
However, on the out-of-domain datasets, we can see that Mistral-TableGPT underperforms the original Table-GPT.
We attribute such performance differences to the differences between the base models.
Since GPT-3.5 is stronger than these open-source 7B models, its innate table understanding ability as well as its generalization ability leads to better performance on these out-of-domain table datasets for Table-GPT.
This reinforces our motivations of conducting the comparisons using the same base model, as \textit{the performance difference may be because of the base model's capability}, therefore we need the same base model to conduct an apple-to-apple comparison.

% \begin{table}[t]
%     \centering
%     \small
%     \renewcommand{\arraystretch}{1.3}
%     \setlength\tabcolsep{2pt}
%     \scalebox{1}{
%     % \resizebox{\linewidth}{!}{
%     \begin{tabular}{c|cc|cccc}
%     \toprule
%    & Beer & SS-DI  & C & SS-CF & WikiTQ & Efth\\
%    & (F1) & (Acc)  & (F1) & (Acc) & (Acc) & (F1)\\
%    \midrule
%     13K & 98.90 & 45.85 & \textbf{29.35} & 21.17 & 29.18 & 66.75\\
%     66K & \textbf{100.00} & \textbf{46.35} & 23.75 & \textbf{25.33} & \textbf{29.77} & \textbf{68.27} \\
%     \bottomrule
%     \end{tabular}
%     }
%     \caption{Performance comparison between training Mistral v0.3 7B Instruct model on the 13K versus 66K instances from \citet{li2023table}.
%     }
%     \label{tab:mistral-tablegpt-side-finding-part}
% \end{table}


\paragraph{Side Findings.}
There is a smaller training set provided by \citet{li2023table} containing 13K training instances.
We report the performance comparison by training the Mistral v0.3 7B Instruct model on the two sets in \Cref{tab:mistral-tablegpt-side-finding-all}
We do not find a significant performance boost when we use the larger 66K dataset.
And on one of the out-of-domain datasets, C, training on 13K instances even yields a better score of 29.4 than training on 66K instances' 23.8.
This echoes with the findings by \citet{zhou2024lima} that limited instruction tuning instances are able to yield a strong model.

\begin{table*}[t]
\small
\centering
\renewcommand{\arraystretch}{1.3}
\setlength\tabcolsep{3pt}
% \scalebox{1}{
\resizebox{\linewidth}{!}{
\begin{tabular}{rccccccccccccccccc}
\hline
\multicolumn{1}{c|}{\multirow{4}{*}{\textbf{\begin{tabular}[c]{@{}c@{}}Train\\Data\end{tabular}}}} & \multicolumn{8}{c|}{\textit{Real}} & \multicolumn{8}{c}{\textit{Synthesized}}\\
\cline{2-17}
& \multicolumn{5}{|c|}{\textbf{Table QA}} & \multicolumn{2}{c|}{\textbf{Fact Veri.}} & \multicolumn{1}{c|}{\textbf{Tab2T}} & \multicolumn{7}{c|}{\textbf{Schema Reasoning}} &  \multicolumn{1}{c}{\textbf{Misc.}}\\
\cline{2-17}
& \multicolumn{1}{|c|}{FeT} & \multicolumn{1}{c|}{HiT}  & \multicolumn{1}{c|}{TabM} & \multicolumn{1}{c|}{TAT}  & \multicolumn{1}{c|}{Wiki}  & \multicolumn{1}{c|}{TabF} & \multicolumn{1}{c|}{Inf} & \multicolumn{1}{c|}{ToT} & \multicolumn{1}{c|}{Beer} & \multicolumn{1}{c|}{DeepM}  & \multicolumn{1}{c|}{DI} & \multicolumn{1}{c|}{ED}  & \multicolumn{1}{c|}{C}  & \multicolumn{1}{c|}{CF} & \multicolumn{1}{c|}{CTA} & \multicolumn{1}{c}{TabB\textsubscript{eval}} \\
 \cline{2-17}
& \multicolumn{1}{|c|}{BLEU} & \multicolumn{1}{c|}{Acc} & \multicolumn{1}{c|}{Acc} & \multicolumn{1}{c|}{Acc} & \multicolumn{1}{c|}{Acc} & \multicolumn{1}{c|}{Acc} & \multicolumn{1}{c|}{Acc} & \multicolumn{1}{c}{BLEU} & \multicolumn{1}{|c|}{F1} & \multicolumn{1}{c|}{Recall} & \multicolumn{1}{c|}{Acc} & \multicolumn{1}{c|}{F1} & \multicolumn{1}{c|}{F1} & \multicolumn{1}{c|}{Acc} & \multicolumn{1}{c|}{F1} & \multicolumn{1}{c}{R-L}\\
\hline
% GPT-3.5$^\dagger$ & 26.49 & 43.62 &  46.3 & 31.38 & 53.13 & 67.41 & 56.00 & 16.81 \\
% GPT-4$^\dagger$ &  21.70 & 48.40 &  67.1 & 55.81 & 68.40 & 74.40 & 58.60 & 12.21 \\
% \hline
% Two decimals
% $\crown$ Base & \textbf{20.04} & 35.54 & 66.94 & 17.98 & 27.88 & \textbf{62.33} & \textbf{42.76} & 11.52 & \textbf{97.25} & 42.86 & 27.90 & 24.06 & \textbf{30.25} & 19.08 & 63.80 & 18.89 \\
% % Table-GPT$^{13K}$ & M-v0.3 7B Instruct & 19.55 & 39.33 & 34.71 & 18.94 & 27.23  & 55.66 & 36.14 & 5.31 & \\
% M-TableLlama & \g 38.66 & \g 70.58 & 71.23 & 5.59 & 23.83 & \g86.83 & 27.74 & \textbf{28.54} & 25.82 & 70.00 & 13.45  & 25.12 & 17.35 & \hspace{0.5em}0.48  & 34.85 & \textbf{19.59}\\
% $\crown$ M-TableLLM & 10.20 & \textbf{44.07} & \textbf{74.99} & 25.02 & 32.30 & 11.91 & 15.39 & 6.70 & 45.05 & \textbf{78.57} & \textbf{33.10} & \textbf{43.12} & 25.55 & 14.98 &  66.92 & 3.67 \\
% M-TableBenchLLM & 7.87 & \textbf{44.07} & 70.62 & \textbf{25.74} & \textbf{37.43} & 36.50 & 27.50 & 3.51 & 88.46 & 50.00 & 32.05 & 20.30 & 27.35 & 13.32 & \textbf{72.22} & \g27.22\\
% M-TableGPT & 19.54 & 35.80 & 62.19 & 14.13 & 25.46 & 61.37 & 35.83 & 4.47 & \g100.00 & \g97.96 & \g46.35 & \g46.02 & 23.75 & \textbf{25.33} & 68.27 & 13.06 \\ 
$\crown$ N/A & \textbf{20.0} & 35.5 & 66.9 & 18.0 & 27.9 & \textbf{62.3} & \textbf{42.8} & 11.5 & \textbf{97.2} & 42.9 & 27.9 & 24.1 & \textbf{30.2} & 19.1 & 63.8 & 18.9 \\
TableLlama & \g 38.7 & \g 70.6 & 71.2 & 5.6 & 23.8 & \g86.8 & 27.7 & \textbf{28.5} & 25.8 & 70.0 & 13.4  & 25.1 & 17.4 & \hspace{0.5em}0.5  & 34.9 & \textbf{19.6}\\
$\crown$ TableLLM & 10.2 & \textbf{44.1} & \textbf{75.0} & 25.0 & 32.3 & 11.9 & 15.4 & 6.7 & 45.0 & \textbf{78.6} & \textbf{33.1} & \textbf{43.1} & 25.6 & 15.0 &  66.9 & 3.7 \\
TableBench & 7.9 & \textbf{44.1} & 70.6 & \textbf{25.7} & \textbf{37.4} & 36.5 & 27.5 & 3.5 & 88.5 & 50.0 & 32.0 & 20.3 & 27.4 & 13.3 & \textbf{72.2} & \g27.2\\
TableGPT & 19.5 & 35.8 & 62.2 & 14.1 & 25.5 & 61.4 & 35.8 & 4.5 & \g100.0 & \g98.0 & \g46.4 & \g46.0 & 23.8 & \textbf{25.3} & 68.3 & 13.1 \\ 
% \hline
% Mistral\textsubscript{Nemo} Instruct (12B)  & 16.25 &  &  &  &  & 54.53 & 56.74 & 13.18 \\
% Mistral\textsubscript{Nemo}-TableLlama & - &  &  &  &  & - & 27.85 & 33.61 \\
% Mistral\textsubscript{Nemo}-TableLLM & 16.54 &   &   &   &  &  56.22 & 58.78 & 11.72 \\
% Mistral\textsubscript{Nemo}-TableBenchLLM &  16.84  &   &   &  &  & 70.66 & 65.70 & 13.11 \\
% Mistral\textsubscript{Nemo}-TableGPT & 28.80  &   &   &  &  & 63.75 & 65.63 & 21.18\\
% \hline
% OLMo Instruct & 6.02 & 27.27 & 54.44 & 14.31 & 19.45 & 38.24 & 21.39 & 5.06\\
% OLMo-TableLlama & \g36.79 & \g67.93 & 72.92 & 9.86 & 6.68 & \g83.75 & 14.96 & 20.82 \\
% OLMo-TableLLM & 9.72 & 35.48 & 65.48 & 17.68 & 26.66 & 40.61 & 16.94 & 8.89\\
% OLMo-TableBenchLLM & 3.80 & 28.34 & 62.63 & 15.57 & 34.02 & 30.94 & 6.48 & 7.50 \\
% OLMo-TableGPT & 9.31 & 27.21 & 65.60 & 14.61 & 16.41 & 44.94 & 32.96 & 11.42 \\
% \hline
% Phi-7B-Instruct & 5.02 & 39.58 & 76.14 & 12.99 & 29.72 & 65.32 & 62.28 & 1.42 \\
% Phi-TableLlama & \g38.06 & \g63.57 & 74.76 & 18.34 & 46.29 & \g86.19 & 54.26 & 29.57 \\
% Phi-TableLLM & 18.19 & 45.27 & 81.23 & 24.11 & 37.68 & 69.64 & 44.63 & 8.08\\
% Phi-TableBenchLLM & 10.03 & 3.47 & 82.97 & 20.51 & 34.62 & 68.01 & 65.26 & 0.90 \\
% Phi-TableGPT & 24.79 & 45.08 & 76.78 & 15.57 & 29.97 & 71.05 & 67.02 & 14.01\\
\hline
\end{tabular}
}
\caption{Out-of-domain evaluation of Mistral v0.3 7B Instruct model fine-tuned with different training data.
``N/A'' denotes the untuned Mistral v0.3 7B Instruct model.
% $\dagger$ denotes performance reported by \citet{zhang-etal-2024-tablellama, gou2023critic, srivastava2024assessing} except for InfoTabs, where we subsample 500 examples to prompt the GPT models ourselves.
The number is in gray if the ``train data'' includes the training set for the corresponding dataset.
``$\crown$'' indicates the training data that leads to the most number of top performance for these table datasets.
Under ``Train Data'', names refer to the datasets used in \Cref{sec: reimpl-compare} for fine-tuning (e.g. TableLlama refers to the training data in \Cref{subsec: reimpl-tablellama}).
% \zshe{why there are two crowns?}
% \dnaihao{TableLLM same source does not lead to advantage.}
}
\label{tab: out-of-domain-mistral}
\end{table*}



\section{Out-of-Domain Table Tasks Evaluation}
\label{sec: real-world-eval}

In \Cref{sec: real-world-eval,sec: general-eval}, we evaluate our fine-tuned models from \Cref{sec: reimpl-compare}.
% We also report the performance of the untuned base model on these datasets.

\subsection{Datasets}
We evaluate these models on eight existing real-world datasets covering the tasks of table question answering (table QA), table fact verification, and table-to-text generation.
\textbf{FeTaQA (FeT)} \citep{nan-etal-2022-fetaqa} is a free-form table QA dataset sourced from Wikipedia-based tables.
\textbf{HiTab (HiT)} \citep{cheng-etal-2022-hitab} is a table QA dataset sourced from statistical reports and Wikipedia pages on hierarchical tables.
\textbf{TabMWP (TabM)} \citep{lu2022dynamic} is an open-domain grade-level table question-answering dataset involving mathematical reasoning.
\textbf{TATQA (TAT)} \citep{zhu-etal-2021-tat} is a table QA dataset sourced from real-world financial reports.
\textbf{WikiTQ (Wiki)} \citep{pasupat-liang-2015-compositional} is a table QA dataset sourced from Wikipedia.
\textbf{TabFact (TabF)} \citep{chen2019tabfact} is a table fact verification dataset sourced from Wikipedia.
\textbf{InfoTabs (Inf)} \citep{gupta-etal-2020-infotabs} is a table fact verification dataset with human-written textual hypotheses based on tables extracted from Wikipedia info-boxes.
\textbf{ToTTo (ToT)} \citep{parikh-etal-2020-totto} is a table-to-text dataset sourced from Wikipedia tables.

In addition, we evaluate these models on eight synthesized datasets including \textbf{Beer}, \textbf{DeepM}, \textbf{Spreadsheet-DI (DI)}, \textbf{Spreadsheet-Real (ED)}, \textbf{Column-No-Separator (C)}, \textbf{Spreadsheet-CF (CF)}, and \textbf{Efthymiou (CTA)} \citep{li2023table} on schema reasoning ability (introduced in \Cref{subsec: mistral-tablegpt}), and \textbf{TabB\textsubscript{eval}} \citep{wu2024tablebenchcomprehensivecomplexbenchmark} on miscellaneous table tasks.
% including table fact verification, numerical reasoning, data analysis, etc.
% Following \citet{wu2024tablebenchcomprehensivecomplexbenchmark}, we report the overall ROUGE-L scores for the test set.
We include a more detailed evaluation setup in \Cref{app-sec: out-real-world-setup} and provide examples from these datasets in \Cref{app-sec: dataset-examples}.

\begin{table}[t]
    \small
    \centering
    \renewcommand{\arraystretch}{1.3}
    \setlength\tabcolsep{2pt}
    \scalebox{1}{
    % \resizebox{\linewidth}{!}{
    \begin{tabular}{cp{6.5cm}}
        \toprule
        Table & \multicolumn{1}{c}{\grayc \begin{tabular}[c]{@{}cccc@{}}Year&Name&Translated Name&Type\\\hline 1961&Hallaç&Carder&Short\\\hline...\end{tabular}} \\
        \midrule
        Q & \grayc How many works did Leyla Erbil publish in total? \\
        \midrule
        \multirow{3}{*}{Gold} & \grayc Leyla Erbil published a total of 11 works. 
        \textit{This can be determined by counting the number of entries in the ``Name'' column in the provided table}.\\
        \bottomrule
    \end{tabular}
    }
    \caption{An example of the training example from TableLLM.
    The reasoning part is in italics. }
    \label{tab:tablellm-example}
\end{table}


% \begin{table}[t]
%     \small
%     \centering
%     \renewcommand{\arraystretch}{1.3}
%     \setlength\tabcolsep{2pt}
%     % \scalebox{1}{
%     \resizebox{\linewidth}{!}{
%     \begin{tabular}{cp{5.5cm}}
%         \toprule
%         \multirow{2}{*}{Gold} & \grayc In his rookie season, Hawk led with 119 total tackles. \\
%         \midrule
%         M-TableLlama & \grayc Hawk played 119 games in his career. \\
%         \midrule
%         \multirow{4}{*}{M-TableLLM} & \grayc The given table cells show the total number of games played by A. J. Hawk, which is 119. The description should state that A. J. Hawk played a total of 119 games.\\
%         \bottomrule
%     \end{tabular}
%     }
%     \caption{An example of the predictions on ToTTo dataset.
%     ``M'' denotes the Mistral model.
%     }
%     \label{tab:totto-example}
% \end{table}



\subsection{Results for Same Base Model, Different Training Data}

\Cref{tab: out-of-domain-mistral} presents our fine-tuned Mistral models' performance on various out-of-domain table datasets.
\Cref{app-sec: more-results} presents the performance of all the models we have fine-tuned in \Cref{sec: reimpl-compare}.


% \patricng{Maybe add this in contribution. In general, I think the paper structure can be reformatted to be based on "RQ" (research question). What is the effect of training on synthetic data on OOD real-world table tasks?}

\paragraph{Base model's performance.}
We find that \textit{the base model is a strong baseline}.
In \Cref{tab: out-of-domain-mistral}, the original Mistral model maintains the best performance on five out-of-domain table datasets.
For the original OLMo and Phi model in \Cref{tab: out-of-domain-complete}, they demonstrate competitive performance compared to the fine-tuned models.
For instance, on the InfoTabs dataset, the untuned Phi model yields 62.3 compared to the best performance of 67.0. 
The original OLMo model achieves 50.5 on the Beer dataset, outperforming the fine-tuned models that are based on OLMo.
This demonstrates that through pre-training and general instruction tuning, these models have acquired innate table understanding ability, which echos with the findings by \citet{li2023table,deng-etal-2024-tables}

\begin{table*}[t]
\small
\centering
\renewcommand{\arraystretch}{1.3}
\setlength\tabcolsep{3pt}
% \scalebox{1}{
\resizebox{\linewidth}{!}{
\begin{tabular}{cccccccccccccccccc}
\hline
\multicolumn{1}{c|}{\multirow{4}{*}{\textbf{\begin{tabular}[c]{@{}c@{}}Base\\Model\end{tabular}}}} & \multicolumn{8}{c|}{\textit{Real}} & \multicolumn{8}{c}{\textit{Synthesized}}\\
\cline{2-17}
& \multicolumn{5}{|c|}{\textbf{Table QA}} & \multicolumn{2}{c|}{\textbf{Fact Veri.}} & \multicolumn{1}{c|}{\textbf{Tab2T}} & \multicolumn{7}{c|}{\textbf{Schema Reasoning}} &  \multicolumn{1}{c}{\textbf{Misc.}}\\
\cline{2-17}
& \multicolumn{1}{|c|}{FeT} & \multicolumn{1}{c|}{HiT}  & \multicolumn{1}{c|}{TabM} & \multicolumn{1}{c|}{TAT}  & \multicolumn{1}{c|}{Wiki}  & \multicolumn{1}{c|}{TabF} & \multicolumn{1}{c|}{Inf} & \multicolumn{1}{c|}{ToT} & \multicolumn{1}{c|}{Beer} & \multicolumn{1}{c|}{DeepM}  & \multicolumn{1}{c|}{DI} & \multicolumn{1}{c|}{ED}  & \multicolumn{1}{c|}{C}  & \multicolumn{1}{c|}{CF} & \multicolumn{1}{c|}{CTA} & \multicolumn{1}{c}{TabB\textsubscript{eval}} \\
 \cline{2-17}
& \multicolumn{1}{|c|}{BLEU} & \multicolumn{1}{c|}{Acc} & \multicolumn{1}{c|}{Acc} & \multicolumn{1}{c|}{Acc} & \multicolumn{1}{c|}{Acc} & \multicolumn{1}{c|}{Acc} & \multicolumn{1}{c|}{Acc} & \multicolumn{1}{c}{BLEU} & \multicolumn{1}{|c|}{F1} & \multicolumn{1}{c|}{Recall} & \multicolumn{1}{c|}{Acc} & \multicolumn{1}{c|}{F1} & \multicolumn{1}{c|}{F1} & \multicolumn{1}{c|}{Acc} & \multicolumn{1}{c|}{F1} & \multicolumn{1}{c}{R-L}\\
\hline
% GPT-3.5$^\dagger$ & 26.49 & 43.62 &  46.3 & 31.38 & 53.13 & 67.41 & 56.00 & 16.81 \\
% GPT-4$^\dagger$ &  21.70 & 48.40 &  67.1 & 55.81 & 68.40 & 74.40 & 58.60 & 12.21 \\
% \hline
% M & 10.20 & 44.07 & 74.99 & \textbf{25.02} & 32.30 & 11.91 & 15.39 & 6.70 & 45.05 & \textbf{78.57} & 33.10 & \textbf{43.12} & 25.55 & 14.98 &  66.92 & 3.67 \\
M & 10.2 & 44.1 & 75.0 & \textbf{25.0} & 32.3 & 11.9 & 15.4 & 6.7 & 45.0 & \textbf{78.6} & 33.1 & \textbf{43.1} & 25.6 & 15.0 &  66.9 & 3.7 \\
% \hline
% Mistral\textsubscript{Nemo} Instruct (12B)  & 16.25 &  &  &  &  & 54.53 & 56.74 & 13.18 \\
% Mistral\textsubscript{Nemo}-TableLlama & - &  &  &  &  & - & 27.85 & 33.61 \\
% Mistral\textsubscript{Nemo}-TableLLM & 16.54 &   &   &   &  &  56.22 & 58.78 & 11.72 \\
% Mistral\textsubscript{Nemo}-TableBenchLLM &  16.84  &   &   &  &  & 70.66 & 65.70 & 13.11 \\
% Mistral\textsubscript{Nemo}-TableGPT & 28.80  &   &   &  &  & 63.75 & 65.63 & 21.18\\

% O & 9.72 & 35.48 & 65.48 & 17.68 & 26.66 & 40.61 & 16.94 & \textbf{8.89} & 16.48 & 42.86 & 33.00 & 37.56 & 13.0 & 18.67 & 43.60 & 6.30\\
% P & \textbf{18.19} & \textbf{45.27} & \textbf{81.23} & 24.11 & \textbf{37.68} & \textbf{69.64} & \textbf{44.63} & 8.08  & \textbf{80.22} & 50.00 & \textbf{34.05} & 41.32 & \textbf{27.95} & \textbf{49.52} & \textbf{70.12} & \textbf{27.23} \\
O & 9.7 & 35.5 & 65.5 & 17.7 & 26.7 & 40.6 & 16.9 & \textbf{8.9} & 16.5 & 42.9 & 33.0 & 37.6 & 13.0 & 18.7 & 43.6 & 6.3\\
P & \textbf{18.2} & \textbf{45.3} & \textbf{81.2} & 24.1 & \textbf{37.7} & \textbf{69.6} & \textbf{44.6} & 8.1  & \textbf{80.2} & 50.0 & \textbf{34.0} & 41.3 & \textbf{27.9} & \textbf{49.5} & \textbf{70.1} & \textbf{27.2} \\

\hline
\end{tabular}
}
\caption{Out-of-domain evaluation for different table tasks.
% $\dagger$ denotes performance reported by \citet{zhang-etal-2024-tablellama, gou2023critic, srivastava2024assessing} except for InfoTabs, where we subsample 500 examples to prompt the GPT models ourselves.
Here the models are all trained on the TableLLM training data.
In terms of the base models, ``M'', ``O'', and ``P'' represent Mistral v0.3 7B Instruct, OLMo 7B Instruct, and Phi 3 Small Instruct (7B), respectively.
We make the number bold if it is the best among the three.
}
\label{tab: model-comparison}
\vspace{-0.5em}
\end{table*}

\paragraph{Table QA tasks.}
TableLLM's training data consistently achieves the best (e.g. on HiTab in \Cref{tab: out-of-domain-mistral}) or competitive performance on table QA tasks across all three base models in \Cref{tab: out-of-domain-complete}.
In contrast, though the recipe for TableLlama's training data contains table QA tasks, models trained with the training data from TableLlama underperform those from TableLLM.
We attribute the effectiveness of TableLLM's training data on the table QA task to that when constructing the data, \citet{zhang2024tablellm} leverage LLMs such as GPT-3.5 to enhance the reasoning process.
For example, in the training instance in \Cref{tab:tablellm-example}, in addition to answering the question, the adjusted gold answer incorporates the reasoning of ``counting the number of entries''.
Such training data teach the base models of the underlying reasoning process to reach to the final answer, therefore benefiting its table QA ability.

\paragraph{Table fact verification tasks.}
When fine-tuned on the Mistral model, TableGPT's training data, while slightly underperforming the base model, exhibits the least performance decay.
In addition, in \Cref{tab: out-of-domain-complete}, fine-tuning on TableGPT's training data achieves the best out-of-domain table fact verification performance for both the OLMo and Phi models.
Interestingly, despite TableBench and TableLlama's training data including table fact verification examples, models trained with these datasets still underperform the base model on the InfoTabs dataset, achieving scores of 27.7 and 27.5 compared to the base model's 42.8 (\Cref{tab: out-of-domain-mistral}). This pattern is consistent across all three base models.
% In \Cref{tab: out-of-domain-complete}, both TableBench and TableLlama's training data fail to improve OLMo's performance on InfoTabs, while for Phi, TableBench's training data improves the performance compared to the base model but is still less effective than TableGPT's training data.
Notably, TableGPT's training data do not explicitly include table fact verification examples, yet they yield the most significant improvements. 
We hypothesize that the key to success lies in the reasoning process rather than the superficial similarity of task formats. 
Although TableLlama's training data include tasks like TabFact, which involve table fact verification, the model may rely on domain-specific patterns rather than the authentic reasoning process to output labels such as ``entail'' or ``refute''.
% We hypothesize that it is the reasoning process rather than the similarity of task format matters to achieve competitive performance on table fact verification.
% Though TableLlama's training data involves TabFact, one of the table fact verification tasks, the model may fail to learn the underlying reasons of why to output ``entail'' or ``refute''.
% In other words, these short labels may fail to supervise the model to learn the authentic reasoning process but some domain-specific patterns.
% On the other hands, TableGPT involves tasks such as error detection and schema matching, and involves step-by-step reasoning.
% These tasks, together with the reasoning, may help models on the table fact verification tasks.
In addition, we highlight that the original TableLlama model, as the previous SOTA model on TabFact, yields 2.85 on InfoTabs as reported by \citet{zheng2024multimodal}.
Our fine-tuned Mistral model, though outperforming the original TableLlama on the TabFact dataset, underperforms the untuned Mistral model on InfoTabs.
\textit{Such results highlight the limitations of the SOTA-chasing works} as even TableLlama and our Mistral-TableLlama achieve competitive performance on TabFact, these models still do not generalize to datasets in the same task category.

\paragraph{Transferability across table tasks.}
As aforementioned, though TableGPT's training data do not explicitly include table fact verification tasks, tasks such as error detection and schema matching appear to positively contribute to table fact verification performance.
Moreover, unlike the original TableLlama model, we do not instruction-tune the Mistral model on the extensive amount of table operation data as described in \Cref{subsec: reimpl-tablellama}.
Nevertheless, fine-tuning the model on tasks like table QA and table fact verification still results in a significant performance boost on DeepM, a schema-matching dataset, with a score of 70.0 compared to the base model's 42.9.
Our finding aligns with \citet{zhang-etal-2024-tablellama}'s observation that training solely on HiTab leads to better performance on certain table operation datasets, suggesting that transferability exists across table tasks.

% On text generation tasks, \textit{the wording difference leads to the poor performance of models trained on synthetic data.}
% For FeTaQA and ToTTo, where the model needs to generate a complete response, models trained on synthetic table data performs poorly.
% We attribute this to the wording difference between the real-world data versus the synthetic training data.
% For instance, \Cref{tab:totto-example} shows a comparison between the generation from Mistral-TableLlama, which is trained on real-world data, and Mistral-TableLLM trained on synthetic data.
% Though the meaning in the two responses are similar, both capturing the important information in the table, Mistral-TableLLM demonstrates the pattern of longer generation and a preference of more complicated structure (e.g. the relative clause including ``which'' versus a plain sentence).  
% Currently, the most wide-adopted approach to evaluate these generations is by BLEU scores \citep{zhang-etal-2024-tablellama,zheng2024multimodal}, such an approach may focus too much on the wording rather than the correctness of the generation.
% On the other hand, \citet{xie-etal-2022-unifiedskg, deng-etal-2024-tables} employ manual effort to check the correctness of the generation, which is hard to scale.
% Our findings invite future effort in building robust and comprehensive evaluation on table-to-text generation.

% \begin{table*}[t]
%     \small
% \centering
% \renewcommand{\arraystretch}{1.3}
% \setlength\tabcolsep{2pt}
% \scalebox{1}{
% \begin{tabular}{ccccccccc}
% \hline
% \multicolumn{1}{c|}{\multirow{2}{*}{\textbf{Method}}}
% % & \multicolumn{1}{c|}{\multirow{3}{*}{\textbf{LLM}}} 
% & \multicolumn{1}{c|}{Beer} & \multicolumn{1}{c|}{DeepM}  & \multicolumn{1}{c|}{SS-DI} & \multicolumn{1}{c|}{SS-Real}  & \multicolumn{1}{c|}{C}  & \multicolumn{1}{c|}{SS-CF} & \multicolumn{1}{c|}{Efth} & \multicolumn{1}{c}{TableBench\textsubscript{eval}}\\
%  \cline{2-9}
% & \multicolumn{1}{|c|}{F1} & \multicolumn{1}{c|}{Recall} & \multicolumn{1}{c|}{Acc} & \multicolumn{1}{c|}{F1} & \multicolumn{1}{c|}{F1} & \multicolumn{1}{c|}{Acc} & \multicolumn{1}{c|}{F1} & \multicolumn{1}{c}{ROUGE-L}\\
% \hline
%     Mistral v0.3 7B Instruct & \textbf{97.25} & 42.86 & 27.90 & 24.06 & \textbf{30.25} & 19.08 & 63.80 & 18.89 \\
% Mistral-TableLlama & 25.82 & 70.00 & 13.45  & 25.12 & 17.35 & \hspace{0.5em}0.48  & 34.85 & \textbf{19.59}\\
% Mistral-TableLLM & 45.05 & \textbf{78.57} & \textbf{33.10} & \textbf{43.12} & 25.55 & 14.98 &  66.92 & 3.67 \\
% Mistral-TableBenchLLM & 88.46 & 50.00 & 32.05 & 20.30 & 27.35 & 13.32 & \textbf{72.22} & \g27.22 \\
% Mistral-TableGPT & \g100.00 & \g97.96 & \g46.35 & \g46.02 & 23.75 & \textbf{25.33} & 68.27 & 13.06 \\
% \hline
% % Mistral\textsubscript{Nemo} Instruct (12B) & 93.96 & 71.43 & 39.50 & 4.21 & 37.40 & 31.81 & 65.82 & 21.76 \\
% % Mistral\textsubscript{Nemo}-TableLlama & 45.05 & 57.14 & 21.50 & 23.09 & 34.00 & 28.00 & 65.57 & 21.66 \\
% % Mistral\textsubscript{Nemo}-TableLLM & 85.71 & 42.86 & 39.55 & 53.80 & 35.45 & 34.01 & 66.33 & 18.46 \\
% % Mistral\textsubscript{Nemo}-TableBenchLLM & 92.86 & 85.71 & 42.65 & 5.59 & 39.55 & 35.08 & 66.67 & - \\
% % Mistral\textsubscript{Nemo}-TableGPT & - & - & - & - & 35.45 & 34.01 & 66.33 & 22.98 \\
% % \hline 
% OLMo & 50.55 & 35.71 & 28.90 & 14.10 & 15.00 & 16.23 & 54.46 & 7.58 \\
% OLMo-TableLlama & 0.00 & 7.14 & 21.20 & 14.64 & 14.75 & 10.70 & 23.48 & 17.06\\
% OLMo-TableLLM & 16.48 & 42.86 & 33.00 & 37.56 & 13.0 & 18.67 & 43.60 & 6.30 \\
% OLMo-TableBenchLLM & 43.40 & 16.59 & 36.60 & 28.64 & 18.15 & 21.22 & 46.55 & \g19.31 \\
% OLMo-TableGPT & \g96.15 & \g100.00 & \g45.40 & 35.34  & 19.95 & 29.31 & 62.46  & 13.66 \\
% \hline
% Phi & 95.05 & 42.86 & 31.90 & 49.70 & 30.65 & 43.40 & 71.55 & 8.30 \\
% Phi-TableLlama & 95.60 & 35.71 & 4.35 & 19.42 & 27.95 & 36.50 & 43.86 & 22.37\\
% Phi-TableLLM & 80.22 & 50.00 & 34.05 & 41.32 & 27.95 & 49.52 & 70.12 & 27.23 \\
% Phi-TableBenchLLM & 95.05 & 28.57 & 35.85 & 53.76 & 31.10 & 46.25 & 76.68 & \g27.76 \\
% Phi-TableGPT & \g98.90 & \g98.81 & \g49.35 & \g55.37 & 24.80 & 45.18 & 68.35 & 26.14 \\
% \hline
%     \end{tabular}
%     }
%     \caption{
%     Out-of-domain evaluation for the synthetic datasets.
%     We put `--' if the model has been in-domain tuned on these datasets.
%     The number is in bold if it is the best.
%     }
%     \label{tab: out-of-domain-synth}
% \end{table*}


\begin{figure*}[t]
    \centering
    \begin{subfigure}[t]{0.9\linewidth}
        \centering
        
\includegraphics[width=\linewidth]{imgs/bar_general_diff/bar_general_legend_m.crop.pdf}
    \end{subfigure}%
    \vspace*{0.4em}
     \begin{subfigure}[t]{0.32\textwidth}
        \centering
        \includegraphics[width=\linewidth]{imgs/bar_general_diff/bar_general_m_modified.pdf}
        \caption{Mistral v0.3 7B Instruct.}
    \end{subfigure}%
    % \hspace{1em}
     \begin{subfigure}[t]{0.32\textwidth}
        \centering
        \includegraphics[width=\linewidth]{imgs/bar_general_diff/bar_general_o_modified.pdf}
        \caption{OLMo 7B Instruct.}
    \end{subfigure}%
    % \hspace{1em}
     \begin{subfigure}[t]{0.32\textwidth}
        \centering
        \includegraphics[width=\linewidth]{imgs/bar_general_diff/bar_general_p_modified.pdf}
        \caption{Phi 3 Small Instruct (7B).}
    \end{subfigure}%

    \caption{Performance of fine-tuned models trained on different data (e.g. TableLlama) on general benchmarks. 
    The green and red hatched bars represent performance gains or losses relative to the base model, respectively.
    % We report the difference between the fine-tuned versus the base model for Mistral v0.3 7B Instruct, OLMo 7B Instruct, and Phi 3 Small Instruct (7B).
    On IFEval, unlike other models, the Mistral model shows a significant performance drop, underscoring the impact of innate model capabilities on preserving general performance after domain-specific fine-tuning.
    % \dnaihao{original capability; select one or two datasets and show; MMLU and IFEval}
    % \dnaihao{use white bar for base model.}
    }
    \label{fig:general-benchmark-performance-difference}
    \vspace{-0.5em}
\end{figure*}

\subsection{Results for Same Training Data, Different Base Models}
\Cref{tab: model-comparison} presents the performance of models fine-tuned from three base models, all on the TableLLM training data.

\paragraph{The effects of the training data depend on the base model.}
In \Cref{tab: model-comparison}, when all fine-tuned on TableLLM's training data, there is a significant performance gap among the three models.
For instance, on WikiTQ, fine-tuning the OLMo model yields 26.7, fine-tuning the Mistral model yields 32.3, while fine-tuning the Phi model yields 37.7.
We attribute the performance difference to the varying innate capabilities of each model.
% In addition, we observe that \textit{a model's performance on general benchmarks can serve as a good indicator to show its capability in table tasks}.
In \Cref{tab:general} in \Cref{app-sec: more-results}, the original Phi model outperforms both the Mistral and OLMo model on four out of five general benchmarks, which aligns with its superior performance on most table tasks after fine-tuning.
% On the other hands, considering that prior table LLMs (introduced in \Cref{tab: tablellm-hyperparameter}) leverage different baseline models, it is hard to disentangle the effectiveness of the training data versus the base model.
Moreover, as switching to stronger base models can lead to better performance even with the same training data, it is possible that existing models may not exhaust the potential of their corresponding training data in their originally reported results.

\paragraph{Strong base model leads to better performance.}
In \Cref{tab: out-of-domain-complete}, the best performance for a single dataset is typically achieved by fine-tuning the base model, which outperforms the other two models when untuned.
For instance, TabMWP's overall best performance is achieved by fine-tuning the Phi model with the TableBench training data, and the original Phi model achieves 76.1, outperforming the original Mistral's 66.9 and the original OLMo's 54.4.
TATQA's overall best performance is achieved by fine-tuning the Mistral model with TableBench training data, and the original Mistral model achieves 18.0, outperforming the original OLMo's 14.3 and the original Phi's 13.0.
This suggests that \textit{practitioners can select the base model by comparing their performance on downstream tasks prior to fine-tuning}, which can save the effort of training all candidate base models before deciding which one to use.

% \dnaihao{Highlight one set of training data may have different performance with different base model.}



\section{General Tasks Evaluation}
\label{sec: general-eval}

In this section, we evaluate these models on general benchmarks to understand how table instruction tuning impacts the models' general capabilities. 
Ideally, a model should maintain its general capabilities as much as possible after the instruction tuning process, because a model that loses significant knowledge during tuning may struggle to serve end-users in real-world applications.

\subsection{Datasets}

\textbf{MMLU} \citep{hendryckstest2021} examines the general ability of the model on 57 tasks including elementary mathematics, US history, computer science, etc.
We adopt the 5-shot setup.
\textbf{MMLU\textsubscript{Pro}} \citep{wang2024mmlupro} is an enhanced benchmark evaluating the general ability of the model, which contains up to ten options and eliminates the trivial questions in MMLU.
We adopt the 5-shot setup.
\textbf{AI2ARC} \citep{clark2018think} is a reasoning benchmark containing natural, grade-school questions. 
We adopt the 0-shot setup and report the accuracy score on the challenging set.
\textbf{GPQA} \citep{rein2023gpqa} is a reasoning benchmark containing questions in biology, physics, and chemistry written by domain experts. 
We adopt a 0-shot setup and report the accuracy score on its main set.
\textbf{IFEval} \citep{zhou2023instruction} is a dataset evaluating the general instruction following ability of the model containing instructions such as ``return the answer in JSON format''.
We report the instance-level strict accuracy defined by \citet{zhou2023instruction}.
We include a more detailed setup in \Cref{app-sec: general-benchmark-eval-setup} for our evaluation process and provide examples from these datasets in \Cref{app-sec: dataset-examples}.

\subsection{Results and Analysis}

\Cref{fig:general-benchmark-performance-difference} presents the results of our models on the general benchmarks.
\Cref{tab:general} in \Cref{app-sec: more-results} presents the complete performance of the base model, our fine-tuned models, and the corresponding difference that we plot in \Cref{fig:general-benchmark-performance-difference}.
We find that on MMLU, MMLU\textsubscript{Pro}, AI2ARC, and GPQA, \textit{our fine-tuned models do not compromise too much of the base models' general capabilities.}
On AI2ARC, the score for Mistral-TableGPT is even slightly higher than the base model.
Such performance improvement is likely due to the fact that many table tasks involve reasoning over tables, which may enhance the model's general reasoning ability.
On IFEval, models fine-tuned from the Mistral model suffer a significant performance drop of over 20 points compared to the original model.
However, models fine-tuned from the Phi model even improve the base model's performance.
We attribute such discrepancy to the difference in the base model's innate characteristics.
For instance, certain base models may be more robust in terms of acquiring new capabilities while maintaining their original capabilities, or the examples of the downstream tasks happen to align well with the examples the model has seen during its general training process.
Our finding suggests that domain-specific tuning does not necessarily lead to performance decay on general benchmarks, and it heavily depends on the base model's innate characteristics.
We provide additional discussion on the effects of model scales for both out-of-domain table tasks evaluation and general tasks evaluation in \Cref{app-sec: more-results}.



\section{Conclusion}
To conduct an apples-to-apples comparison, we train three 7B Instruct models based on the table instruction tuning data proposed in the prior works.
As a side product, we achieve the new SOTA performance on HiTab.
We are the first to decouple the factors of training data versus base models and provide analysis on each side.
In addition, we conduct evaluations on the general benchmarks to investigate how domain-specific fine-tuning may influence the model's general capabilities.
We hope our work provides future directions for research on structured data.


\section*{Limitations}

We believe our work presents a comprehensive evaluation over a diverse set of table benchmarks and the general benchmarks.
In addition, we want to stress the massive training effort we have invested in, as the training data provided in the existing works can be as large as 100K.
As a side product, we have achieved the new SOTA performance on HiTab dataset, and provide the first open-source model replication of existing closed-source table LLMs such as Table-GPT.
However, there exists other datasets proposed by the researchers which can be further used to evaluate these models' capabilities, and by no means we can exhaust all of them in this paper.
We encourage future efforts in comprehensively evaluating these table LLMs' capabilities, and we believe our work has laid a solid foundation to decoupling the contributions of training data and base models, and further enhancing our understanding of table instruction tuning.

% \section*{Impact Statements}
% Tables represent a structured and systematic way of organizing vast amounts of information, serving as the backbone for numerous applications such as medical diagnostics \citep{akhtar-etal-2022-pubhealthtab}, virtual personal assistants, and customer relationship management systems \citep{hemphill1990atis, dahl1994expanding}. 
% Over the past decades, the community has introduced numerous table-related datasets to address diverse scenarios involving tabular data, including table question-answering datasets such as WikiTQA \citep{pasupat-liang-2015-compositional}, table-to-text generation datasets such as ToTTo \citep{parikh-etal-2020-totto}, and table fact verification datasets such as TabFact \citep{chen2019tabfact}.
% To tackle these tasks, researchers have proposed various specialized table models \citep{yin-etal-2020-tabert, herzig-etal-2020-tapas, yang-etal-2022-tableformer, xie-etal-2022-unifiedskg}. 
% The diversity and impact of these datasets and models underscore the growing focus on tabular data and its wide-ranging applications.

% In this work, we focus on isolating the contributions of training data used by existing table-specific LLMs by training consistent base models and benchmarking their performance. 
% The models that we evaluate include the Mistral v0.3 7B Instruct model \citep{jiang2023mistral}, OLMo 7B Instruct \citep{groeneveld2024olmo}, and Phi 3 Small Instruct (7B) \citep{abdin2024phi}, with additional experiments on the Phi 3 Mini Instruct (4B) detailed in \Cref{app-sec: more-results}.

% Our replicated models serve as strong baselines for future research on structured data.
% We provide a comprehensive evaluation of both their performance on table-specific tasks and their trade-offs with general capabilities. 
% Our findings offer actionable insights into which training data contributes most significantly to table tasks and propose strategies for efficiently constructing LLMs optimized for tabular data. 

% Our work lays a foundation for further exploration into table instruction tuning and its applications, enabling more robust and effective models for structured data tasks.

% Unlike unstructured text, tables systematically organize large amounts of information, serving as the foundation for numerous applications, including medical medical diagnostics \citep{akhtar-etal-2022-pubhealthtab}, virtual personal assistants, and customer relationship management systems \citep{hemphill1990atis, dahl1994expanding}.
% In the past decades, the community has proposed numerous table-related datasets to explore different scenarios involving tables, including table question-answering datasets such as WikiTQA \citep{pasupat-liang-2015-compositional}, table-to-text generation datasets such as ToTTo \citep{parikh-etal-2020-totto}, table fact verification datasets such as TabFact \citep{chen2019tabfact}. 
% To tackle these tasks, researchers have been putting a significant amount of effort into building strong table models \citep{yin-etal-2020-tabert, herzig-etal-2020-tapas, yang-etal-2022-tableformer, xie-etal-2022-unifiedskg}.
% Such diversity and the high impact of these table-related datasets and models underscore the wide application scenarios for models specialized in table tasks.

\section*{Ethical Considerations}
In this work, we isolate the contributions of training data proposed by the existing table LLMs by training the same base models and comparing their performance.
The base models we have used in this work include Mistral v0.3 7B Instruct model \citep{jiang2023mistral}, OLMo 7B Instruct \citep{groeneveld2024olmo}, and Phi 3 Small Instruct (7B) \citep{abdin2024phi}.
We conduct additional studies on Phi 3 Mini Instruct (4B) in \Cref{app-sec: more-results}.
Foundational models like Mistral v0.3 7B Instruct model are susceptible to jail-breaking instructions \citep{wei2024jailbroken} and may lead to harmful behaviors.
Our objective in this work is to understand the limitations of the existing table instruction tuning, and we urge practitioners to stick to the good purpose when developing or using our models.
Our replicated models can serve as baseline models for future research on structured data, and we provide a holistic evaluation of these models on both table tasks and how they compromise their general capabilities.
Our results lead to various findings on what training data helps the models most on these table tasks, and how to construct LLMs specialized in tables efficiently.



\bibliography{aclanthology,custom}
% \bibliographystyle{iclr2024_conference}

\newpage
\appendix
\newpage

\section{Experimental Setup}
\label{app-sec: experimental-setup}

We run our experiments on 1 server node with 8 A100, each with 48 GB GPU memory.
We set the batch size to 16 in our training process.


\section{Out-of-Domain Evaluation Setup}
\label{app-sec: out-real-world-setup}
For FeTaQA, we use the BLEU4 score following \citet{nan-etal-2022-fetaqa}.
For ToTTo, we follow \citet{xie-etal-2022-unifiedskg} to report the BLEU4 scores over multiple references.
We adopt the evaluation script from the original HiTab, TabMWP, TATQA, and WikiTQ repository on GitHub.
For these table QA tasks, we notice that since the fine-tuned models may not follow instructions such as ``generate in the JSON format'', we do not pose any constraints to these models in terms of the generation format.
Instead, we use Haiku 3.5\footnote{\url{https://www.anthropic.com/claude/haiku}} to extract the answer entity from the model generation. 
For TabFact and InfoTabs, we report the accuracy by checking if only the gold answer appears in the prediction.

% To avoid any false positives, following \citet{zheng2024multimodal}, we ask the model to return the final result as JSON in the format \{``answer'': ``\textless YOUR ANSWER\textgreater''\} for QA tasks HiTab, TabMWP, TATQA, and WikiTQ, and fact verification tasks TabFact and InfoTabs.

In terms of the test set format, we use the exact same test set for FeTaQA, HiTab, TATQA, and ToTTo as \citet{zhang-etal-2024-tablellama} with the Markdown table format.
For TabMWP, WikiTQ, and InfoTabs, etc., we follow the original data format.
Specifically, TabMWP uses `$|$' to separate columns, and WikiTQ and InfoTabs use HTML format to represent tables.

\section{General Evaluation Setup}
\label{app-sec: general-benchmark-eval-setup}

For MMLU, MMLU\textsubscript{Pro}, AI2ARC, and GPQA, as they are all multi-choice question-answering datasets, our objective is to select the most appropriate completion among a set of given options based on the provided context. 
Following \citet{touvron2023llama}, we select the completion with the highest likelihood given the provided context. 
As we evaluate the model based on their selection of the letter choice of ``A'', ``B'', etc., we do not normalize the likelihood by the number of characters in the completion.


\section{More Results}
\label{app-sec: more-results}


\subsection{Out-of-domain Table Tasks Evaluation}

% \begin{figure}[t]
% \centering
% \begin{subfigure}
%         \centering
%         \includegraphics[width=\linewidth]{imgs/spider_plots/performance_spider_legend.crop.pdf}
%     \end{subfigure}%
    
%     \vspace*{1mm}
    
%  \begin{subfigure}
%         \centering
%         \includegraphics[width=\linewidth]{imgs/spider_plots/performance_spider.crop.pdf}
%     \end{subfigure}%
%         \caption{
%         Performance comparison between ``M'', the Mistral v0.3 7B Instruct model, and Mistral-TableLlama, Mistral-TableLLM, Mistral-TableBenchLLM, and Mistral-TableGPT.
%         ``Real'', ``Synth'' denote the out-of-domain real and synthetic table tasks' performance, ``General'' denotes the performance on the general benchmarks.
%         }
%          \label{fig:performance-spider}
% \end{figure}


% In \Cref{fig:performance-spider}, for out-of-domain real-world table tasks, we calculate the model's average performance on TabMWP, TATQA, WikiTQ, InfoTabs, and ToTTo introduced in \Cref{sec: real-world-eval}.
% For out-of-domain synthetic table tasks, we calculate the model's average performance on C, SS-CF, and Efth introduced in \Cref{sec: synth-eval}.
% We select these datasets as they are out-of-domain for every model we train.
% For general tasks, we calculate the model's average performance on MMLU, MMLU\textsubscript{Pro}, AI2ARC, GPQA, and IFEval introduced in \Cref{sec: general-eval}.

\begin{table*}[t]
\small
\centering
\renewcommand{\arraystretch}{1.3}
\setlength\tabcolsep{2pt}
% \scalebox{1}{
\resizebox{\linewidth}{!}{
\begin{tabular}{rccccccccccccccccc}
\hline
\multicolumn{1}{c|}{\multirow{4}{*}{\textbf{\begin{tabular}[c]{@{}c@{}}Train\\Data\end{tabular}}}} & \multicolumn{8}{c|}{\textit{Real}} & \multicolumn{8}{c}{\textit{Synthesized}}\\
\cline{2-17}
& \multicolumn{5}{|c|}{\textbf{Table QA}} & \multicolumn{2}{c|}{\textbf{Fact Veri.}} & \multicolumn{1}{c|}{\textbf{Tab2Text}} & \multicolumn{7}{c|}{\textbf{Schema Reasoning}} &  \multicolumn{1}{c}{\textbf{Misc.}}\\
\cline{2-17}
& \multicolumn{1}{|c|}{FeT} & \multicolumn{1}{c|}{HiT}  & \multicolumn{1}{c|}{TabM} & \multicolumn{1}{c|}{TAT}  & \multicolumn{1}{c|}{Wiki}  & \multicolumn{1}{c|}{TabF} & \multicolumn{1}{c|}{Inf} & \multicolumn{1}{c|}{ToT} & \multicolumn{1}{c|}{Beer} & \multicolumn{1}{c|}{DeepM}  & \multicolumn{1}{c|}{DI} & \multicolumn{1}{c|}{ED}  & \multicolumn{1}{c|}{C}  & \multicolumn{1}{c|}{CF} & \multicolumn{1}{c|}{CTA} & \multicolumn{1}{c}{TabB\textsubscript{eval}} \\
 \cline{2-17}
& \multicolumn{1}{|c|}{BLEU} & \multicolumn{1}{c|}{Acc} & \multicolumn{1}{c|}{Acc} & \multicolumn{1}{c|}{Acc} & \multicolumn{1}{c|}{Acc} & \multicolumn{1}{c|}{Acc} & \multicolumn{1}{c|}{Acc} & \multicolumn{1}{c}{BLEU} & \multicolumn{1}{|c|}{F1} & \multicolumn{1}{c|}{Recall} & \multicolumn{1}{c|}{Acc} & \multicolumn{1}{c|}{F1} & \multicolumn{1}{c|}{F1} & \multicolumn{1}{c|}{Acc} & \multicolumn{1}{c|}{F1} & \multicolumn{1}{c}{ROUGE-L}\\
\hline
% GPT-3.5$^\dagger$ & 26.49 & 43.62 &  46.3 & 31.38 & 53.13 & 67.41 & 56.00 & 16.81 \\
% GPT-4$^\dagger$ &  21.70 & 48.40 &  67.1 & 55.81 & 68.40 & 74.40 & 58.60 & 12.21 \\
% \hline
% $\crown$ M-Base & \textbf{20.04} & 35.54 & 66.94 & 17.98 & 27.88 & \textbf{62.33} & \textbf{42.76} & 11.52 & \textbf{97.25} & 42.86 & 27.90 & 24.06 & \textbf{30.25} & 19.08 & 63.80 & 18.89 \\
% % Table-GPT$^{13K}$ & M-v0.3 7B Instruct & 19.55 & 39.33 & 34.71 & 18.94 & 27.23  & 55.66 & 36.14 & 5.31 & \\
% M-TableLlama & \g 38.66 & \g 70.58 & 71.23 & 5.59 & 23.83 & \g86.83 & 27.74 & \textbf{28.54} & 25.82 & 70.00 & 13.45  & 25.12 & 17.35 & \hspace{0.5em}0.48  & 34.85 & \textbf{19.59}\\
% $\crown$ M-TableLLM & 10.20 & \textbf{44.07} & \textbf{74.99} & 25.02 & 32.30 & 11.91 & 15.39 & 6.70 & 45.05 & \textbf{78.57} & \textbf{33.10} & \textbf{43.12} & 25.55 & 14.98 &  66.92 & 3.67 \\
% M-TableBenchLLM & 7.87 & \textbf{44.07} & 70.62 & \textbf{25.74} & \textbf{37.43} & 36.50 & 27.50 & 3.51 & 88.46 & 50.00 & 32.05 & 20.30 & 27.35 & 13.32 & \textbf{72.22} & \g27.22\\
% M-TableGPT & 19.54 & 35.80 & 62.19 & 14.13 & 25.46 & 61.37 & 35.83 & 4.47 & \g100.00 & \g97.96 & \g46.35 & \g46.02 & 23.75 & \textbf{25.33} & 68.27 & 13.06 \\ 
% % \hline
% % Mistral\textsubscript{Nemo} Instruct (12B)  & 16.25 &  &  &  &  & 54.53 & 56.74 & 13.18 \\
% % Mistral\textsubscript{Nemo}-TableLlama & - &  &  &  &  & - & 27.85 & 33.61 \\
% % Mistral\textsubscript{Nemo}-TableLLM & 16.54 &   &   &   &  &  56.22 & 58.78 & 11.72 \\
% % Mistral\textsubscript{Nemo}-TableBenchLLM &  16.84  &   &   &  &  & 70.66 & 65.70 & 13.11 \\
% % Mistral\textsubscript{Nemo}-TableGPT & 28.80  &   &   &  &  & 63.75 & 65.63 & 21.18\\
% \hline
% O-Base & 6.02 & 27.27 & 54.44 & 14.31 & 19.45 & 38.24 & 21.39 & 5.06 & \textbf{50.55} & 35.71 & 28.90 & 14.10 & 15.00 & 16.23 & 54.46 & 7.58\\
% O-TableLlama & \g36.79 & \g67.93 & \textbf{72.92} & 9.86 & 6.68 & \g83.75 & 14.96 & \textbf{20.82} & 0.00 & 7.14 & 21.20 & 14.64 & 14.75 & 10.70 & 23.48 & \textbf{17.06}\\
% $\crown$ O-TableLLM & \textbf{9.72} & \textbf{35.48} & 65.48 & \textbf{17.68} & 26.66 & 40.61 & 16.94 & 8.89 & 16.48 & \textbf{42.86} & 33.00 & \textbf{37.56} & 13.0 & 18.67 & 43.60 & 6.30\\
% O-TableBenchLLM & 3.80 & 28.34 & 62.63 & 15.57 & \textbf{34.02} & 30.94 & 6.48 & 7.50 & 43.40 & 16.59 & \textbf{36.60} & 28.64 & 18.15 & 21.22 & 46.55 & \g19.31 \\
% $\crown$ O-TableGPT & 9.31 & 27.21 & 65.60 & 14.61 & 16.41 & \textbf{44.94} & \textbf{32.96} & 11.42 & \g96.15 & \g100.00 & \g45.40 & 35.34  & \textbf{19.95} & \textbf{29.31} & \textbf{62.46}  & 13.66 \\
% \hline 
% P-Base & 5.02 & 39.58 & 76.14 & 12.99 & 29.72 & 65.32 & 62.28 & 1.42  & 95.05 & 42.86 & 31.90 & 49.70 & 30.65 & 43.40 & \textbf{71.55} & 8.30 \\
% P-TableLlama & \g38.06 & \g63.57 & 74.76 & 18.34 & \textbf{46.29} & \g86.19 & 54.26 & \textbf{29.57}  & \textbf{95.60} & 35.71 & 4.35 & 19.42 & 27.95 & 36.50 & 43.86 & 22.37\\
% $\crown$ P-TableLLM & 18.19 & \textbf{45.27} & 81.23 & \textbf{24.11} & 37.68 & 69.64 & 44.63 & 8.08  & 80.22 & \textbf{50.00} & 34.05 & 41.32 & 27.95 & \textbf{49.52} & 70.12 & \textbf{27.23} \\
% P-TableBenchLLM & 10.03 & 3.47 & \textbf{82.97} & 20.51 & 34.62 & 68.01 & 65.26 & 0.90 & 95.05 & 28.57 & \textbf{35.85} & \textbf{53.76} & \textbf{31.10} & 46.25 & 76.68 & \g27.76\\
% P-TableGPT & \textbf{24.79} & 45.08 & 76.78 & 15.57 & 29.97 & \textbf{71.05} & \textbf{67.02} & 14.01 & \g98.90 & \g98.81 & \g49.35 & \g55.37 & 24.80 & 45.18 & 68.35 & 26.14 \\
\multicolumn{17}{l}{\grayc Mistral v0.3 7B Instruct} \\
$\crown$ N/A & \textbf{20.0} & 35.5 & 66.9 & 18.0 & 27.9 & \textbf{62.3} & \textbf{42.8} & 11.5 & \red\textbf{97.2} & 42.9 & 27.9 & 24.1 & \textbf{30.2} & 19.1 & 63.8 & 18.9 \\
% Table-GPT$^{13K}$ & M-v0.3 7B Instruct & 19.6 & 39.3 & 34.7 & 18.9 & 27.2  & 55.7 & 36.1 & 5.3 & \\
TableLlama & \g 38.7 & \g 70.6 & 71.2 & 5.6 & 23.8 & \g86.8 & 27.7 & \textbf{28.5} & 25.8 & 70.0 & 13.4  & 25.1 & 17.4 & \hspace{0.5em}0.5  & 34.9 & \textbf{19.6}\\
$\crown$ TableLLM & 10.2 & \textbf{44.1} & \textbf{75.0} & 25.0 & 32.3 & 11.9 & 15.4 & 6.7 & 45.0 & \red\textbf{78.6} & \textbf{33.1} & \textbf{43.1} & 25.6 & 15.0 &  66.9 & 3.7 \\
TableBenchLLM & 7.9 & \textbf{44.1} & 70.6 & \red\textbf{25.7} & \textbf{37.4} & 36.5 & 27.5 & 3.5 & 88.5 & 50.0 & 32.0 & 20.3 & 27.4 & 13.3 & \textbf{72.2} & \g27.2\\
TableGPT & 19.5 & 35.8 & 62.2 & 14.1 & 25.5 & 61.4 & 35.8 & 4.5 & \g100.0 & \g98.0 & \g46.4 & \g46.0 & 23.8 & \textbf{25.3} & 68.3 & 13.1 \\ 
\hline
\multicolumn{17}{l}{\grayc OLMo 7B Instruct} \\
N/A & 6.0 & 27.3 & 54.4 & 14.3 & 19.4 & 38.2 & 21.4 & 5.1 & \textbf{50.5} & 35.7 & 28.9 & 14.1 & 15.0 & 16.2 & 54.5 & 7.6\\
TableLlama & \g36.8 & \g67.9 & \textbf{72.9} & 9.9 & 6.7 & \g83.8 & 15.0 & \textbf{20.8} & 0.0 & 7.1 & 21.2 & 14.6 & 14.8 & 10.7 & 23.5 & \textbf{17.1}\\
$\crown$ TableLLM & \textbf{9.7} & \textbf{35.5} & 65.5 & \textbf{17.7} & 26.7 & 40.6 & 16.9 & 8.9 & 16.5 & \textbf{42.9} & 33.0 & \textbf{37.6} & 13.0 & 18.7 & 43.6 & 6.3\\
TableBenchLLM & 3.8 & 28.3 & 62.6 & 15.6 & \textbf{34.0} & 30.9 & 6.5 & 7.5 & 43.4 & 16.6 & \red\textbf{36.6} & 28.6 & 18.1 & 21.2 & 46.5 & \g19.3 \\
$\crown$ TableGPT & 9.3 & 27.2 & 65.6 & 14.6 & 16.4 & \textbf{44.9} & \textbf{33.0} & 11.4 & \g96.2 & \g100.0 & \g45.4 & 35.3  & \textbf{19.9} & \textbf{29.3} & \textbf{62.5}  & 13.7 \\
\hline 
\multicolumn{17}{l}{\grayc Phi 3 Small Instruct (7B)} \\
N/A & 5.0 & 39.6 & 76.1 & 13.0 & 29.7 & 65.3 & 62.3 & 1.4  & 95.0 & 42.9 & 31.9 & 49.7 & 30.6 & 43.4 & 71.5 & 8.3 \\
TableLlama & \g38.1 & \g63.6 & 74.8 & 18.3 & \red\textbf{46.3} & \g86.2 & 54.3 & \red\textbf{29.6}  & \textbf{95.6} & 35.7 & 4.3 & 19.4 & 27.9 & 36.5 & 43.9 & 22.4\\
$\crown$ TableLLM & 18.2 & \red\textbf{45.3} & 81.2 & \textbf{24.1} & 37.7 & 69.6 & 44.6 & 8.1  & 80.2 & \textbf{50.0} & 34.0 & 41.3 & 27.9 & \red\textbf{49.5} & 70.1 & \red\textbf{27.2} \\
$\crown$ TableBenchLLM & 10.0 & 3.5 & \red\textbf{83.0} & 20.5 & 34.6 & 68.0 & 65.3 & 0.9 & 95.0 & 28.6 & \textbf{35.9} & \red\textbf{53.8} & \red\textbf{31.1} & 46.2 & \red\textbf{76.7} & \g27.8\\
TableGPT & \red\textbf{24.8} & 45.1 & 76.8 & 15.6 & 30.0 & \red\textbf{71.0} & \red\textbf{67.0} & 14.0 & \g98.9 & \g98.8 & \g49.4 & \g55.4 & 24.8 & 45.2 & 68.3 & 26.1 \\
\hline
\end{tabular}
}
\caption{Out-of-domain evaluation for different table tasks.
% $\dagger$ denotes performance reported by \citet{zhang-etal-2024-tablellama, gou2023critic, srivastava2024assessing} except for InfoTabs, where we subsample 500 examples to prompt the GPT models ourselves.
The number is in gray if the model's training data contains the training set corresponding to the dataset.
We make the number bold if it is the best among the same model, and the number red if it is the best across all the models.
Mistral v0.3 7B Instruct, OLMo 7B Instruct, and Phi 3 Small Instruct (7B) indicate the base model on which we apply the training data, respectively.
``$\crown$'' indicates the model has the most number of top performance across all the datasets with respect to the same base model.
}
\label{tab: out-of-domain-complete}
\end{table*}

\paragraph{Comparison across 7B models.}
\Cref{tab: out-of-domain-complete} presents performance scores for our fine-tuned models from \Cref{sec: reimpl-compare} and their corresponding base models.
We find that when training with different base models in \Cref{tab: out-of-domain-complete}, TableLLM's training data consistently yield the best performance on the most out-of-domain table datasets.


% \begin{table*}[t]
% \small
% \centering
% \renewcommand{\arraystretch}{1.3}
% \setlength\tabcolsep{2pt}
% % \scalebox{1}{
% \resizebox{\linewidth}{!}{
% \begin{tabular}{cccccccccccccccccc}
% \hline
% \multicolumn{1}{c|}{\multirow{4}{*}{\textbf{Method}}} & \multicolumn{8}{c|}{\textit{Real}} & \multicolumn{8}{c}{\textit{Synthesized}}\\
% \cline{2-17}
% & \multicolumn{5}{|c|}{\textbf{Table QA}} & \multicolumn{2}{c|}{\textbf{Fact Veri.}} & \multicolumn{1}{c|}{\textbf{Tab2Text}} & \multicolumn{7}{c|}{\textbf{Schema Reasoning}} &  \multicolumn{1}{c}{\textbf{Misc.}}\\
% \cline{2-17}
% & \multicolumn{1}{|c|}{FeT} & \multicolumn{1}{c|}{HiT}  & \multicolumn{1}{c|}{TabM} & \multicolumn{1}{c|}{TAT}  & \multicolumn{1}{c|}{Wiki}  & \multicolumn{1}{c|}{TabF} & \multicolumn{1}{c|}{Inf} & \multicolumn{1}{c|}{ToT} & \multicolumn{1}{c|}{Beer} & \multicolumn{1}{c|}{DeepM}  & \multicolumn{1}{c|}{DI} & \multicolumn{1}{c|}{ED}  & \multicolumn{1}{c|}{C}  & \multicolumn{1}{c|}{CF} & \multicolumn{1}{c|}{CTA} & \multicolumn{1}{c}{TabB\textsubscript{eval}} \\
%  \cline{2-17}
% & \multicolumn{1}{|c|}{BLEU} & \multicolumn{1}{c|}{Acc} & \multicolumn{1}{c|}{Acc} & \multicolumn{1}{c|}{Acc} & \multicolumn{1}{c|}{Acc} & \multicolumn{1}{c|}{Acc} & \multicolumn{1}{c|}{Acc} & \multicolumn{1}{c}{BLEU} & \multicolumn{1}{|c|}{F1} & \multicolumn{1}{c|}{Recall} & \multicolumn{1}{c|}{Acc} & \multicolumn{1}{c|}{F1} & \multicolumn{1}{c|}{F1} & \multicolumn{1}{c|}{Acc} & \multicolumn{1}{c|}{F1} & \multicolumn{1}{c}{ROUGE-L}\\
% \hline
% % GPT-3.5$^\dagger$ & 26.49 & 43.62 &  46.3 & 31.38 & 53.13 & 67.41 & 56.00 & 16.81 \\
% % GPT-4$^\dagger$ &  21.70 & 48.40 &  67.1 & 55.81 & 68.40 & 74.40 & 58.60 & 12.21 \\
% % \hline
% % $\crown$ M-Base & \textbf{20.04} & 35.54 & 66.94 & 17.98 & 27.88 & \textbf{62.33} & \textbf{42.76} & 11.52 & \textbf{97.25} & 42.86 & 27.90 & 24.06 & \textbf{30.25} & 19.08 & 63.80 & 18.89 \\
% % % Table-GPT$^{13K}$ & M-v0.3 7B Instruct & 19.55 & 39.33 & 34.71 & 18.94 & 27.23  & 55.66 & 36.14 & 5.31 & \\
% % M-TableLlama & \g 38.66 & \g 70.58 & 71.23 & 5.59 & 23.83 & \g86.83 & 27.74 & \textbf{28.54} & 25.82 & 70.00 & 13.45  & 25.12 & 17.35 & \hspace{0.5em}0.48  & 34.85 & \textbf{19.59}\\
% % $\crown$ M-TableLLM & 10.20 & \textbf{44.07} & \textbf{74.99} & 25.02 & 32.30 & 11.91 & 15.39 & 6.70 & 45.05 & \textbf{78.57} & \textbf{33.10} & \textbf{43.12} & 25.55 & 14.98 &  66.92 & 3.67 \\
% % M-TableBenchLLM & 7.87 & \textbf{44.07} & 70.62 & \textbf{25.74} & \textbf{37.43} & 36.50 & 27.50 & 3.51 & 88.46 & 50.00 & 32.05 & 20.30 & 27.35 & 13.32 & \textbf{72.22} & \g27.22\\
% % M-TableGPT & 19.54 & 35.80 & 62.19 & 14.13 & 25.46 & 61.37 & 35.83 & 4.47 & \g100.00 & \g97.96 & \g46.35 & \g46.02 & 23.75 & \textbf{25.33} & 68.27 & 13.06 \\ 
% % % \hline
% % % Mistral\textsubscript{Nemo} Instruct (12B)  & 16.25 &  &  &  &  & 54.53 & 56.74 & 13.18 \\
% % % Mistral\textsubscript{Nemo}-TableLlama & - &  &  &  &  & - & 27.85 & 33.61 \\
% % % Mistral\textsubscript{Nemo}-TableLLM & 16.54 &   &   &   &  &  56.22 & 58.78 & 11.72 \\
% % % Mistral\textsubscript{Nemo}-TableBenchLLM &  16.84  &   &   &  &  & 70.66 & 65.70 & 13.11 \\
% % % Mistral\textsubscript{Nemo}-TableGPT & 28.80  &   &   &  &  & 63.75 & 65.63 & 21.18\\
% % \hline
% % O-Base & 6.02 & 27.27 & 54.44 & 14.31 & 19.45 & 38.24 & 21.39 & 5.06 & \textbf{50.55} & 35.71 & 28.90 & 14.10 & 15.00 & 16.23 & 54.46 & 7.58\\
% % O-TableLlama & \g36.79 & \g67.93 & \textbf{72.92} & 9.86 & 6.68 & \g83.75 & 14.96 & \textbf{20.82} & 0.00 & 7.14 & 21.20 & 14.64 & 14.75 & 10.70 & 23.48 & \textbf{17.06}\\
% % $\crown$ O-TableLLM & \textbf{9.72} & \textbf{35.48} & 65.48 & \textbf{17.68} & 26.66 & 40.61 & 16.94 & 8.89 & 16.48 & \textbf{42.86} & 33.00 & \textbf{37.56} & 13.0 & 18.67 & 43.60 & 6.30\\
% % O-TableBenchLLM & 3.80 & 28.34 & 62.63 & 15.57 & \textbf{34.02} & 30.94 & 6.48 & 7.50 & 43.40 & 16.59 & \textbf{36.60} & 28.64 & 18.15 & 21.22 & 46.55 & \g19.31 \\
% % $\crown$ O-TableGPT & 9.31 & 27.21 & 65.60 & 14.61 & 16.41 & \textbf{44.94} & \textbf{32.96} & 11.42 & \g96.15 & \g100.00 & \g45.40 & 35.34  & \textbf{19.95} & \textbf{29.31} & \textbf{62.46}  & 13.66 \\
% % \hline 
% % P-Base & 5.02 & 39.58 & 76.14 & 12.99 & 29.72 & 65.32 & 62.28 & 1.42  & 95.05 & 42.86 & 31.90 & 49.70 & 30.65 & 43.40 & \textbf{71.55} & 8.30 \\
% % P-TableLlama & \g38.06 & \g63.57 & 74.76 & 18.34 & \textbf{46.29} & \g86.19 & 54.26 & \textbf{29.57}  & \textbf{95.60} & 35.71 & 4.35 & 19.42 & 27.95 & 36.50 & 43.86 & 22.37\\
% % $\crown$ P-TableLLM & 18.19 & \textbf{45.27} & 81.23 & \textbf{24.11} & 37.68 & 69.64 & 44.63 & 8.08  & 80.22 & \textbf{50.00} & 34.05 & 41.32 & 27.95 & \textbf{49.52} & 70.12 & \textbf{27.23} \\
% % P-TableBenchLLM & 10.03 & 3.47 & \textbf{82.97} & 20.51 & 34.62 & 68.01 & 65.26 & 0.90 & 95.05 & 28.57 & \textbf{35.85} & \textbf{53.76} & \textbf{31.10} & 46.25 & 76.68 & \g27.76\\
% % P-TableGPT & \textbf{24.79} & 45.08 & 76.78 & 15.57 & 29.97 & \textbf{71.05} & \textbf{67.02} & 14.01 & \g98.90 & \g98.81 & \g49.35 & \g55.37 & 24.80 & 45.18 & 68.35 & 26.14 \\
% P-base & 10.26 & 33.96 & 81.46 & 13.11 & 26.40 & 41.80 & 55.52 & 8.31 & 93.96 & 53.17 & 33.80 & 26.83 & 24.20 & 18.79 & 65.82 & 16.11 \\
% P-TableLlama & \g 35.12 & \g 48.55 & 72.79 & 17.86 & 34.81 & \g 81.62 & 47.74 & 17.07 & 91.76 & 62.70 & 20.35 & 7.50 & 23.00 & 18.73 & 63.13 & 21.29 \\
% P-TableLLM & 11.58 & 38.26 & 80.74 & 22.67 & 31.56 & 46.15 & 42.61 & 7.77 & 92.31 & 64.29 & 36.80 & 44.71 & 24.20 & 20.93 & 63.80 & 19.22 \\
% P-TableBenchLLM &  10.66 & 37.37 & 85.30 & 20.63 & 36.83 & 48.68 & 56.54 & 8.31 & 90.11 & 38.89 & 35.75 & 29.22 & 23.80 & 21.11 & 66.58 &\g25.60\\
% P-TableGPT & 16.43 & 31.06 & 70.06 & 16.18 & 22.67 & 53.24 & 50.06 & 12.36 & \g 74.18 & \g 56.35 & \g 43.35 & 38.02 & 22.80 & 22.83 & 65.57 & 21.16 \\
% \hline
% P-Base & 5.02 & 39.58 & 76.14 & 12.99 & 29.72 & 65.32 & 62.28 & 1.42  & 95.05 & 42.86 & 31.90 & 49.70 & 30.65 & 43.40 & \textbf{71.55} & 8.30 \\
% P-TableLlama & \g38.06 & \g63.57 & 74.76 & 18.34 & \textbf{46.29} & \g86.19 & 54.26 & \textbf{29.57}  & \textbf{95.60} & 35.71 & 4.35 & 19.42 & 27.95 & 36.50 & 43.86 & 22.37\\
% $\crown$ P-TableLLM & 18.19 & \textbf{45.27} & 81.23 & \textbf{24.11} & 37.68 & 69.64 & 44.63 & 8.08  & 80.22 & \textbf{50.00} & 34.05 & 41.32 & 27.95 & \textbf{49.52} & 70.12 & \textbf{27.23} \\
% P-TableBenchLLM & 10.03 & 3.47 & \textbf{82.97} & 20.51 & 34.62 & 68.01 & 65.26 & 0.90 & 95.05 & 28.57 & \textbf{35.85} & \textbf{53.76} & \textbf{31.10} & 46.25 & 76.68 & \g27.76\\
% P-TableGPT & \textbf{24.79} & 45.08 & 76.78 & 15.57 & 29.97 & \textbf{71.05} & \textbf{67.02} & 14.01 & \g98.90 & \g98.81 & \g49.35 & \g55.37 & 24.80 & 45.18 & 68.35 & 26.14 \\
% \hline
% \end{tabular}
% }
% \caption{Out-of-domain evaluation for different table tasks.
% % $\dagger$ denotes performance reported by \citet{zhang-etal-2024-tablellama, gou2023critic, srivastava2024assessing} except for InfoTabs, where we subsample 500 examples to prompt the GPT models ourselves.
% The number is in gray if the model has been in-domain tuned on the dataset.
% We make the number bold if it is the best among the same model, the number red if it is the best across all the models.
% ``M-'', ``O-'', and ``P-'' represent Mistral v0.3 7B Instruct, OLMo 7B Instruct, Phi 3 Small Instruct (7B), respectively.
% ``$\crown$'' indicates the model has the most number of top performance across all the datasets with respect to the same base model.
% }
% \label{tab: out-of-domain-different-sizes}
% \end{table*}

\paragraph{Comparison across different model sizes.}
\Cref{fig:table-benchmark-size-1,fig:table-benchmark-size-2} provide performance comparison between Phi 3 Mini Instruct (4B) versus Phi 3 Small Instruct (7B).
We find that the larger sized model often leads to better performance for both the original model and the model after training on the same set of data.



\begin{table}[t]
    \small
    \centering
    \renewcommand{\arraystretch}{1.3}
    \setlength\tabcolsep{2pt}
    % \scalebox{1}{
    \resizebox{\linewidth}{!}{
    \begin{tabular}{cccccc}
    \hline
        \multicolumn{1}{c|}{\multirow{2}{*}{\textbf{Method}}} & \multicolumn{1}{c|}{MMLU} &  \multicolumn{1}{c|}{MMLU\textsubscript{Pro}} & \multicolumn{1}{c|}{AI2ARC} & \multicolumn{1}{c|}{GPQA} & \multicolumn{1}{c}{IFEval} \\
        \cline{2-6}
        & \multicolumn{1}{|c|}{Acc} & \multicolumn{1}{c|}{Acc} & \multicolumn{1}{c|}{Acc} & \multicolumn{1}{c|}{Acc} & \multicolumn{1}{c}{Acc}\\
        \hline
%      M    & 61.22 & 31.36 & 73.29 & 28.57 & 58.75\\
%      M-TableLlama & 59.36 & 29.48 & 69.62 & 23.66 & 38.01 \\
%      $\Delta$ & \hspace{-0.2em}$\downarrow$ 1.86  & \hspace{-0.2em}$\downarrow$ 1.88& \hspace{-0.2em}$\downarrow$ 3.37 & \hspace{-0.2em}$\downarrow$ 4.91 & \hspace{-0.3em}$\downarrow$20.74\\
%      \midrule
%      M-TableLLM & 61.44 & 29.32 & 74.23 & 25.89 & 29.62\\
%      $\Delta$ & \hspace{-0.2em}$\uparrow$ 0.22 & \hspace{-0.2em}$\downarrow$ 2.04 & \hspace{-0.2em}$\uparrow$ 0.94 & \hspace{-0.2em}$\downarrow$ 2.68 & \hspace{-0.3em}$\downarrow$29.13\\
%      \midrule
%      M-TableBenchLLM & 61.96 & 31.01 & 73.63 & 28.12 & 31.77 \\
%      $\Delta$ & \hspace{-0.2em}$\uparrow$ 0.74 & \hspace{-0.2em}$\downarrow$ 0.35 & \hspace{-0.2em}$\uparrow$ 0.34 & \hspace{-0.2em}$\downarrow$ 0.45 & \hspace{-0.3em}$\downarrow$26.98\\
%      \midrule
%      M-TableGPT & 61.27 & 31.27 & 74.57 & 26.12 & 31.41\\
%      $\Delta$ & \hspace{-0.2em}$\uparrow$ 0.05 & \hspace{-0.2em}$\downarrow$ 0.09 & \hspace{-0.2em}$\uparrow$ 1.28 & \hspace{-0.2em}$\downarrow$ 2.45 & \hspace{-0.3em}$\downarrow$27.34\\
%      \hline
%      % O & 52.56 & 22.54 & 67.58 & 27.90 & 45.56 \\
%      % O-TableLlama & 53.66 & 23.13 & 66.21 & 29.69 & 46.76\\
%      % O-TableLLM & 53.25 & 22.30 & 65.96 & 29.02 & 42.81 \\
%      % O-TableBenchLLM & 53.06 & 21.88 & 67.66 & 28.57 & 45.20 \\
%      % O-TableGPT & 52.86 & 21.90 & 66.81 & 28.79 & 48.92\\
%      % \hline
%      % M\textsubscript{Nemo} & 45.20 & 32.75 & 60.75 & 33.48 & 68.11\\
%      % M\textsubscript{Nemo}-TableLlama &  48.01 & 32.73 & 53.75 & 28.35 & 65.95\\
%      % M\textsubscript{Nemo}-TableLLM & 49.69 & 32.93 & 69.62 & 31.25 & 67.75\\
%      % M\textsubscript{Nemo}-TableBenchLLM & 47.96 & 33.14 & 61.35 & 32.81 & 64.15 \\
%      % M\textsubscript{Nemo}-TableGPT & 49.17 & 32.29 & 56.57 & 31.47 & 68.11\\
%      % \hline
%      % P & 75.67 & 41.23 & 73.12 & 31.03 & 60.67 \\
%      % P-TableLlama & 75.51 & 45.11 & 73.46  & 31.47 & 70.14 \\
%      % P-TableLLM & 75.02 & 42.57 & 73.12 & 30.36 & 64.75 \\
%      % P-TableBenchLLM & 75.72 & 43.26 & 60.75 & 28.79 & 63.31\\
%      % P-TableGPT & 75.14 & 40.08 & 72.61 & 32.37 & 70.02\\
%      % \hline
%      \hline
% O & 52.56 & 22.54 & 67.58 & 27.90 & 45.56 \\
% O-TableLlama & 53.66 & 23.13 & 66.21 & 29.69 & 46.76 \\
% $\Delta$ & $\uparrow$ 1.10 & $\uparrow$ 0.59 & $\downarrow$ 1.37 & $\uparrow$ 1.99 & $\uparrow$ 1.20 \\\hline
% O-TableLLM & 53.25 & 22.30 & 65.96 & 29.02 & 42.81 \\
% $\Delta$ & $\uparrow$ 0.69 & $\downarrow$ 0.34 & $\downarrow$ 1.62 & $\uparrow$ 1.92 & $\downarrow$ 2.75 \\\hline
% O-TableBenchLLM & 53.06 & 21.88 & 67.66 & 28.57 & 45.20 \\
% $\Delta$ & $\uparrow$ 0.50 & $\downarrow$ 0.66 & $\uparrow$ 0.08 & $\uparrow$ 0.87 & $\downarrow$ 0.36 \\\hline
% O-TableGPT & 52.86 & 21.90 & 66.81 & 28.79 & 48.92 \\
% $\Delta$ & $\uparrow$ 0.30 & $\downarrow$ 0.64 & $\downarrow$ 0.77 & $\uparrow$ 0.79 & $\uparrow$ 3.36 \\
% \hline
% \hline
% P & 75.67 & 41.23 & 73.12 & 31.03 & 60.67 \\
% P-TableLlama & 75.51 & 45.11 & 73.46 & 31.47 & 70.14 \\
% $\Delta$ & $\downarrow$ 0.16 & $\uparrow$ 3.88 & $\uparrow$ 0.34 & $\uparrow$ 0.37 & $\uparrow$ 9.89 \\\hline
% P-TableLLM & 75.02 & 42.57 & 73.12 & 30.36 & 64.75 \\
% $\Delta$ & $\downarrow$ 0.65 & $\uparrow$ 1.34 & $\uparrow$ 0.00 & $\downarrow$ 0.75 & $\uparrow$ 4.08 \\\hline
% P-TableBenchLLM & 75.72 & 43.26 & 60.75 & 28.79 & 63.31 \\
% $\Delta$ & $\uparrow$ 0.05 & $\uparrow$ 2.03 & $\downarrow$ 1.52 & $\downarrow$ 2.14 & $\uparrow$ 2.64 \\\hline
% P-TableGPT & 75.14 & 40.08 & 72.61 & 32.37 & 70.02 \\
% $\Delta$ & $\downarrow$ 0.53 & $\downarrow$ 1.23 & $\downarrow$ 0.34 & $\uparrow$ 1.37 & $\uparrow$ 9.35 \\
% \hline
M    & 61.2 & 31.4 & 73.3 & 28.6 & 58.8 \\
M-TableLlama & 59.4 & 29.5 & 69.6 & 23.7 & 38.0 \\
$\Delta$ & \hspace{-0.2em}$\downarrow$ 1.9 & \hspace{-0.2em}$\downarrow$ 1.9 & \hspace{-0.2em}$\downarrow$ 3.4 & \hspace{-0.2em}$\downarrow$ 4.9 & \hspace{-0.3em}$\downarrow$ 20.7 \\
\midrule
M-TableLLM & 61.4 & 29.3 & 74.2 & 25.9 & 29.6 \\
$\Delta$ & \hspace{-0.2em}$\uparrow$ 0.2 & \hspace{-0.2em}$\downarrow$ 2.0 & \hspace{-0.2em}$\uparrow$ 0.9 & \hspace{-0.2em}$\downarrow$ 2.7 & \hspace{-0.3em}$\downarrow$ 29.1 \\
\midrule
M-TableBenchLLM & 62.0 & 31.0 & 73.6 & 28.1 & 31.8 \\
$\Delta$ & \hspace{-0.2em}$\uparrow$ 0.7 & \hspace{-0.2em}$\downarrow$ 0.4 & \hspace{-0.2em}$\uparrow$ 0.3 & \hspace{-0.2em}$\downarrow$ 0.5 & \hspace{-0.3em}$\downarrow$ 27.0 \\
\midrule
M-TableGPT & 61.3 & 31.3 & 74.6 & 26.1 & 31.4 \\
$\Delta$ & \hspace{-0.2em}$\uparrow$ 0.1 & \hspace{-0.2em}$\downarrow$ 0.1 & \hspace{-0.2em}$\uparrow$ 1.3 & \hspace{-0.2em}$\downarrow$ 2.4 & \hspace{-0.3em}$\downarrow$ 27.3 \\
\hline
\hline
O & 52.6 & 22.5 & 67.6 & 27.9 & 45.6 \\
O-TableLlama & 53.7 & 23.1 & 66.2 & 29.7 & 46.8 \\
$\Delta$ & $\uparrow$ 1.1 & $\uparrow$ 0.6 & $\downarrow$ 1.4 & $\uparrow$ 2.0 & $\uparrow$ 1.2 \\\hline
O-TableLLM & 53.3 & 22.3 & 66.0 & 29.0 & 42.8 \\
$\Delta$ & $\uparrow$ 0.7 & $\downarrow$ 0.3 & $\downarrow$ 1.6 & $\uparrow$ 1.9 & $\downarrow$ 2.8 \\\hline
O-TableBenchLLM & 53.1 & 21.9 & 67.7 & 28.6 & 45.2 \\
$\Delta$ & $\uparrow$ 0.5 & $\downarrow$ 0.7 & $\uparrow$ 0.1 & $\uparrow$ 0.9 & $\downarrow$ 0.4 \\\hline
O-TableGPT & 52.9 & 21.9 & 66.8 & 28.8 & 48.9 \\
$\Delta$ & $\uparrow$ 0.3 & $\downarrow$ 0.6 & $\downarrow$ 0.8 & $\uparrow$ 0.8 & $\uparrow$ 3.4 \\
\hline
\hline
P & 75.7 & 41.2 & 73.1 & 31.0 & 60.7 \\
P-TableLlama & 75.5 & 45.1 & 73.5 & 31.5 & 70.1 \\
$\Delta$ & $\downarrow$ 0.2 & $\uparrow$ 3.9 & $\uparrow$ 0.3 & $\uparrow$ 0.4 & $\uparrow$ 9.9 \\\hline
P-TableLLM & 75.0 & 42.6 & 73.1 & 30.4 & 64.8 \\
$\Delta$ & $\downarrow$ 0.7 & $\uparrow$ 1.3 & $\uparrow$ 0.0 & $\downarrow$ 0.8 & $\uparrow$ 4.1 \\\hline
P-TableBenchLLM & 75.7 & 43.3 & 60.8 & 28.8 & 63.3 \\
$\Delta$ & $\uparrow$ 0.0 & $\uparrow$ 2.0 & $\downarrow$ 1.5 & $\downarrow$ 2.1 & $\uparrow$ 2.6 \\\hline
P-TableGPT & 75.1 & 40.1 & 72.6 & 32.4 & 70.0 \\
$\Delta$ & $\downarrow$ 0.5 & $\downarrow$ 1.2 & $\downarrow$ 0.3 & $\uparrow$ 1.4 & $\uparrow$ 9.4 \\
\hline

    \end{tabular}
    }
    \caption{Evaluation of the models on general benchmarks.
    ``M-'', ``O-'', and ``P-'' represent Mistral v0.3 7B Instruct, OLMo 7B Instruct, Phi 3 Small Instruct (7B), respectively.
    ``$\Delta$'' denotes the performance difference between the instruction-tuned model and its base model.
    }
    \label{tab:general}
\end{table}

\subsection{General Tasks Evaluation}
\paragraph{Comparison across 7B models.}
\Cref{tab:general} presents the performance of the base model, our fine-tuned models, and the corresponding performance difference (also plotted in \Cref{fig:general-benchmark-performance-difference}).
We find that the original model's capability typically decides the fine-tuned models' capability.
With proper training, the original models' capability can be largely preserved even after domain-specific fine-tuning.



% \begin{table}[t]
%     \small
%     \centering
%     \renewcommand{\arraystretch}{1.3}
%     \setlength\tabcolsep{2pt}
%     % \scalebox{1}{
%     \resizebox{\linewidth}{!}{
%     \begin{tabular}{cccccc}
%     \hline
%         \multicolumn{1}{c|}{\multirow{2}{*}{\textbf{Method}}} & \multicolumn{1}{c|}{MMLU} &  \multicolumn{1}{c|}{MMLU\textsubscript{Pro}} & \multicolumn{1}{c|}{AI2ARC} & \multicolumn{1}{c|}{GPQA} & \multicolumn{1}{c}{IFEval} \\
%         \cline{2-6}
%         & \multicolumn{1}{|c|}{Acc} & \multicolumn{1}{c|}{Acc} & \multicolumn{1}{c|}{Acc} & \multicolumn{1}{c|}{Acc} & \multicolumn{1}{c}{Acc}\\
%         \hline
%      M    & 61.22 & 31.36 & 73.29 & 28.57 & 58.75\\
%      M-TableLlama & 59.36 & 29.48 & 69.62 & 23.66 & 38.01 \\
%      $\Delta$ & \hspace{-0.2em}$\downarrow$ 1.86  & \hspace{-0.2em}$\downarrow$ 1.88& \hspace{-0.2em}$\downarrow$ 3.37 & \hspace{-0.2em}$\downarrow$ 4.91 & \hspace{-0.3em}$\downarrow$20.74\\
%      \midrule
%      M-TableLLM & 61.44 & 29.32 & 74.23 & 25.89 & 29.62\\
%      $\Delta$ & \hspace{-0.2em}$\uparrow$ 0.22 & \hspace{-0.2em}$\downarrow$ 2.04 & \hspace{-0.2em}$\uparrow$ 0.94 & \hspace{-0.2em}$\downarrow$ 2.68 & \hspace{-0.3em}$\downarrow$29.13\\
%      \midrule
%      M-TableBenchLLM & 61.96 & 31.01 & 73.63 & 28.12 & 31.77 \\
%      $\Delta$ & \hspace{-0.2em}$\uparrow$ 0.74 & \hspace{-0.2em}$\downarrow$ 0.35 & \hspace{-0.2em}$\uparrow$ 0.34 & \hspace{-0.2em}$\downarrow$ 0.45 & \hspace{-0.3em}$\downarrow$26.98\\
%      \midrule
%      M-TableGPT & 61.27 & 31.27 & 74.57 & 26.12 & 31.41\\
%      $\Delta$ & \hspace{-0.2em}$\uparrow$ 0.05 & \hspace{-0.2em}$\downarrow$ 0.09 & \hspace{-0.2em}$\uparrow$ 1.28 & \hspace{-0.2em}$\downarrow$ 2.45 & \hspace{-0.3em}$\downarrow$27.34\\
%      \hline
%      % O & 52.56 & 22.54 & 67.58 & 27.90 & 45.56 \\
%      % O-TableLlama & 53.66 & 23.13 & 66.21 & 29.69 & 46.76\\
%      % O-TableLLM & 53.25 & 22.30 & 65.96 & 29.02 & 42.81 \\
%      % O-TableBenchLLM & 53.06 & 21.88 & 67.66 & 28.57 & 45.20 \\
%      % O-TableGPT & 52.86 & 21.90 & 66.81 & 28.79 & 48.92\\
%      % \hline
%      % M\textsubscript{Nemo} & 45.20 & 32.75 & 60.75 & 33.48 & 68.11\\
%      % M\textsubscript{Nemo}-TableLlama &  48.01 & 32.73 & 53.75 & 28.35 & 65.95\\
%      % M\textsubscript{Nemo}-TableLLM & 49.69 & 32.93 & 69.62 & 31.25 & 67.75\\
%      % M\textsubscript{Nemo}-TableBenchLLM & 47.96 & 33.14 & 61.35 & 32.81 & 64.15 \\
%      % M\textsubscript{Nemo}-TableGPT & 49.17 & 32.29 & 56.57 & 31.47 & 68.11\\
%      % \hline
%      % P & 75.67 & 41.23 & 73.12 & 31.03 & 60.67 \\
%      % P-TableLlama & 75.51 & 45.11 & 73.46  & 31.47 & 70.14 \\
%      % P-TableLLM & 75.02 & 42.57 & 73.12 & 30.36 & 64.75 \\
%      % P-TableBenchLLM & 75.72 & 43.26 & 60.75 & 28.79 & 63.31\\
%      % P-TableGPT & 75.14 & 40.08 & 72.61 & 32.37 & 70.02\\
%      % \hline
%      \hline
% O & 52.56 & 22.54 & 67.58 & 27.90 & 45.56 \\
% O-TableLlama & 53.66 & 23.13 & 66.21 & 29.69 & 46.76 \\
% $\Delta$ & $\uparrow$ 1.10 & $\uparrow$ 0.59 & $\downarrow$ 1.37 & $\uparrow$ 1.99 & $\uparrow$ 1.20 \\\hline
% O-TableLLM & 53.25 & 22.30 & 65.96 & 29.02 & 42.81 \\
% $\Delta$ & $\uparrow$ 0.69 & $\downarrow$ 0.34 & $\downarrow$ 1.62 & $\uparrow$ 1.92 & $\downarrow$ 2.75 \\\hline
% O-TableBenchLLM & 53.06 & 21.88 & 67.66 & 28.57 & 45.20 \\
% $\Delta$ & $\uparrow$ 0.50 & $\downarrow$ 0.66 & $\uparrow$ 0.08 & $\uparrow$ 0.87 & $\downarrow$ 0.36 \\\hline
% O-TableGPT & 52.86 & 21.90 & 66.81 & 28.79 & 48.92 \\
% $\Delta$ & $\uparrow$ 0.30 & $\downarrow$ 0.64 & $\downarrow$ 0.77 & $\uparrow$ 0.79 & $\uparrow$ 3.36 \\
% \hline
% \hline
% P & 75.67 & 41.23 & 73.12 & 31.03 & 60.67 \\
% P-TableLlama & 75.51 & 45.11 & 73.46 & 31.47 & 70.14 \\
% $\Delta$ & $\downarrow$ 0.16 & $\uparrow$ 3.88 & $\uparrow$ 0.34 & $\uparrow$ 0.37 & $\uparrow$ 9.89 \\\hline
% P-TableLLM & 75.02 & 42.57 & 73.12 & 30.36 & 64.75 \\
% $\Delta$ & $\downarrow$ 0.65 & $\uparrow$ 1.34 & $\uparrow$ 0.00 & $\downarrow$ 0.75 & $\uparrow$ 4.08 \\\hline
% P-TableBenchLLM & 75.72 & 43.26 & 60.75 & 28.79 & 63.31 \\
% $\Delta$ & $\uparrow$ 0.05 & $\uparrow$ 2.03 & $\downarrow$ 1.52 & $\downarrow$ 2.14 & $\uparrow$ 2.64 \\\hline
% P-TableGPT & 75.14 & 40.08 & 72.61 & 32.37 & 70.02 \\
% $\Delta$ & $\downarrow$ 0.53 & $\downarrow$ 1.23 & $\downarrow$ 0.34 & $\uparrow$ 1.37 & $\uparrow$ 9.35 \\
% \hline
% M    & 61.2 & 31.4 & 73.3 & 28.6 & 58.8 \\
% M-TableLlama & 59.4 & 29.5 & 69.6 & 23.7 & 38.0 \\
% $\Delta$ & \hspace{-0.2em}$\downarrow$ 1.9 & \hspace{-0.2em}$\downarrow$ 1.9 & \hspace{-0.2em}$\downarrow$ 3.4 & \hspace{-0.2em}$\downarrow$ 4.9 & \hspace{-0.3em}$\downarrow$ 20.7 \\
% \midrule
% M-TableLLM & 61.4 & 29.3 & 74.2 & 25.9 & 29.6 \\
% $\Delta$ & \hspace{-0.2em}$\uparrow$ 0.2 & \hspace{-0.2em}$\downarrow$ 2.0 & \hspace{-0.2em}$\uparrow$ 0.9 & \hspace{-0.2em}$\downarrow$ 2.7 & \hspace{-0.3em}$\downarrow$ 29.1 \\
% \midrule
% M-TableBenchLLM & 62.0 & 31.0 & 73.6 & 28.1 & 31.8 \\
% $\Delta$ & \hspace{-0.2em}$\uparrow$ 0.7 & \hspace{-0.2em}$\downarrow$ 0.4 & \hspace{-0.2em}$\uparrow$ 0.3 & \hspace{-0.2em}$\downarrow$ 0.5 & \hspace{-0.3em}$\downarrow$ 27.0 \\
% \midrule
% M-TableGPT & 61.3 & 31.3 & 74.6 & 26.1 & 31.4 \\
% $\Delta$ & \hspace{-0.2em}$\uparrow$ 0.1 & \hspace{-0.2em}$\downarrow$ 0.1 & \hspace{-0.2em}$\uparrow$ 1.3 & \hspace{-0.2em}$\downarrow$ 2.4 & \hspace{-0.3em}$\downarrow$ 27.3 \\
% \hline
% \hline
% O & 52.6 & 22.5 & 67.6 & 27.9 & 45.6 \\
% O-TableLlama & 53.7 & 23.1 & 66.2 & 29.7 & 46.8 \\
% $\Delta$ & $\uparrow$ 1.1 & $\uparrow$ 0.6 & $\downarrow$ 1.4 & $\uparrow$ 2.0 & $\uparrow$ 1.2 \\\hline
% O-TableLLM & 53.3 & 22.3 & 66.0 & 29.0 & 42.8 \\
% $\Delta$ & $\uparrow$ 0.7 & $\downarrow$ 0.3 & $\downarrow$ 1.6 & $\uparrow$ 1.9 & $\downarrow$ 2.8 \\\hline
% O-TableBenchLLM & 53.1 & 21.9 & 67.7 & 28.6 & 45.2 \\
% $\Delta$ & $\uparrow$ 0.5 & $\downarrow$ 0.7 & $\uparrow$ 0.1 & $\uparrow$ 0.9 & $\downarrow$ 0.4 \\\hline
% O-TableGPT & 52.9 & 21.9 & 66.8 & 28.8 & 48.9 \\
% $\Delta$ & $\uparrow$ 0.3 & $\downarrow$ 0.6 & $\downarrow$ 0.8 & $\uparrow$ 0.8 & $\uparrow$ 3.4 \\
% \hline
% P & 70.55 & 40.41 & 85.24 & 31.03 & 61.87\\
% P-TableLlama & 70.75 & 40.02 & 85.41 & 31.92 & 56.35 \\
% P-TableLLM & 70.23 & 39.71 & 85.07 & 31.70 & 56.59 \\
% P-TableBenchLLM & 70.42 & 39.82 & 84.81 & 33.26 & 58.87 \\
% P-TableGPT & 70.36 & 39.70 & 85.07 & 31.92 & 60.91 \\
% \hline
% P & 75.67 & 41.23 & 73.12 & 31.03 & 60.67 \\
% P-TableLlama & 75.51 & 45.11 & 73.46 & 31.47 & 70.14 \\
% P-TableLLM & 75.02 & 42.57 & 73.12 & 30.36 & 64.75 \\
% P-TableBenchLLM & 75.72 & 43.26 & 60.75 & 28.79 & 63.31 \\
% P-TableGPT & 75.14 & 40.08 & 72.61 & 32.37 & 70.02 \\
% \hline

%     \end{tabular}
%     }
%     \caption{Evaluation of the models on general benchmarks.
%     ``M-'', ``O-'', and ``P-'' represent Mistral v0.3 7B Instruct, OLMo 7B Instruct, Phi 3 Small Instruct (7B), respectively.
%     ``$\Delta$'' denotes the performance difference between the instruction-tuned model and its base model.
%     }
%     \label{tab:general-different-sizes}
% \end{table}

\begin{figure*}[t]
    \centering
    \begin{subfigure}[t]{0.3\textwidth}
        \centering
        
\includegraphics[width=\linewidth]{imgs/bar_general_diff_size/size_comparison_general_legend.crop.pdf}
    \end{subfigure}%
    
    \vspace*{1em}
     \begin{subfigure}[t]{0.95\textwidth}
        \centering
        \includegraphics[width=\linewidth]{imgs/bar_table_diff_size/size_comparison_tables_0.crop.pdf}
        \caption{No training data, the original model.}
    \end{subfigure}%
    \vspace*{1em}
     \begin{subfigure}[t]{0.95\textwidth}
        \centering
        \includegraphics[width=\linewidth]{imgs/bar_table_diff_size/size_comparison_tables_1.crop.pdf}
        \caption{Training data for TableLlama.}
    \end{subfigure}%
    \vspace*{1em}
     \begin{subfigure}[t]{0.95\textwidth}
        \centering
        \includegraphics[width=\linewidth]{imgs/bar_table_diff_size/size_comparison_tables_2.crop.pdf}
        \caption{Training data for TableLLM.}
    \end{subfigure}%

    \caption{Performance of Phi 3 Mini Instruct (4B) versus Phi 3 Small Instruct (7B) model on different table tasks with different training data.
    % We report the difference between the fine-tuned versus the base model for Mistral v0.3 7B Instruct, OLMo 7B Instruct, and Phi 3 Small Instruct (7B).
    In most cases, the 7B model outperforms the 4B model.
    }
    \label{fig:table-benchmark-size-1}
    
\end{figure*}


\begin{figure*}[t]
    \centering
    \begin{subfigure}[t]{0.3\textwidth}
        \centering
        
\includegraphics[width=\linewidth]{imgs/bar_general_diff_size/size_comparison_general_legend.crop.pdf}
    \end{subfigure}%
    
    \vspace*{1em}
     \begin{subfigure}[t]{0.95\textwidth}
        \centering
        \includegraphics[width=\linewidth]{imgs/bar_table_diff_size/size_comparison_tables_3.crop.pdf}
        \caption{Training data for TableBench.}
    \end{subfigure}%
    \vspace*{1em}
     \begin{subfigure}[t]{0.95\textwidth}
        \centering
        \includegraphics[width=\linewidth]{imgs/bar_table_diff_size/size_comparison_tables_4.crop.pdf}
        \caption{Training data for TableGPT.}
    \end{subfigure}%

    \caption{Performance of Phi 3 Mini Instruct (4B) versus Phi 3 Small Instruct (7B) model on different table tasks with different training data.
    % We report the difference between the fine-tuned versus the base model for Mistral v0.3 7B Instruct, OLMo 7B Instruct, and Phi 3 Small Instruct (7B).
    In most cases, the 7B model outperforms the 4B model.
    }
    \label{fig:table-benchmark-size-2}
    
\end{figure*}

\begin{figure*}[t]
    \centering
    \begin{subfigure}[t]{0.3\textwidth}
        \centering
        
\includegraphics[width=\linewidth]{imgs/bar_general_diff_size/size_comparison_general_legend.crop.pdf}
    \end{subfigure}%
    
    \vspace*{1em}
     \begin{subfigure}[t]{0.30\textwidth}
        \centering
        \includegraphics[width=\linewidth]{imgs/bar_general_diff_size/size_comparisons_general_MMLU.crop.pdf}
        \caption{MMLU.}
    \end{subfigure}%
    \hspace{0.1em}
     \begin{subfigure}[t]{0.30\textwidth}
        \centering
        \includegraphics[width=\linewidth]{imgs/bar_general_diff_size/size_comparisons_general_MMLU_Pro.crop.pdf}
        \caption{MMLU\textsubscript{Pro}.}
    \end{subfigure}%
    \hspace{0.1em}
     \begin{subfigure}[t]{0.30\textwidth}
        \centering
        \includegraphics[width=\linewidth]{imgs/bar_general_diff_size/size_comparisons_general_AI2ARC.crop.pdf}
        \caption{AI2ARC.}
    \end{subfigure}%
    \hspace{0.1em}
     \begin{subfigure}[t]{0.30\textwidth}
        \centering
        \includegraphics[width=\linewidth]{imgs/bar_general_diff_size/size_comparisons_general_GPQA.crop.pdf}
        \caption{GPQA.}
    \end{subfigure}%
    \hspace{0.1em}
     \begin{subfigure}[t]{0.30\textwidth}
        \centering
        \includegraphics[width=\linewidth]{imgs/bar_general_diff_size/size_comparisons_general_IFEval.crop.pdf}
        \caption{IFEval.}
    \end{subfigure}%

    \caption{Performance difference between Phi 3 Mini Instruct (4B) versus Phi 3 Small Instruct (7B) model. 
    On MMLU, MMLU\textsubscript{Pro}, IFEval, the Small (7B) version yields better performance both before and after fine-tuning.
    On GPQA, the two models perform comparably.
    On AI2ARC, the Mini (4B) version yields better performance.
    % We report the difference between the fine-tuned versus the base model for Mistral v0.3 7B Instruct, OLMo 7B Instruct, and Phi 3 Small Instruct (7B).
    }
    \label{fig:general-benchmark-size}
    
\end{figure*}

\paragraph{Comparison across different model sizes.}
\Cref{fig:general-benchmark-size} provides the performance comparison between the Phi 3 Mini Instruct (4B) versus the Phi 3 Small Instruct (7B) model on the five general benchmarks.
On most datasets, the 7B model outperforms the 4B model.
However, on AI2ARC, the 4B model performs better, and on GPQA, the two models perform comparably.
We note that on AI2ARC, we adopt a zero-shot setup, where we do not provide any examples to the models.
The 7B model in this case may not prefer to answer their question at the very beginning, leading to an incorrect probability distribution over the four choices.
For GPQA, as the task itself is challenging, both the 4B and 7B models cannot answer most of them, leading to a comparable performance.

\section{Dataset Examples}
\label{app-sec: dataset-examples}

\subsection{FeTaQA}

\textbf{Input:}
\begin{lstlisting}
[TLE] The Wikipedia page title of this table is Gerhard Bigalk. The Wikipedia section title of this table is Ships attacked. [TAB] | Date | Name | Nationality | Tonnage (GRT) | Fate | [SEP] | 14 June 1941 | St. Lindsay | United Kingdom | 5,370 | Sunk | [SEP] | 21 December 1941 | HMS Audacity | Royal Navy | 11,000 | Sunk | [SEP] | 2 February 1942 | Corilla | Netherlands | 8,096 | Damaged | [SEP] | 4 February 1942 | Silveray | United Kingdom | 4,535 | Sunk | [SEP] | 7 February 1942 | Empire Sun | United Kingdom | 6,952 | Sunk | [SEP] | 16 May 1942 | Nicarao | United States | 1,445 | Sunk | [SEP] | 19 May 1942 | Isabela | United States | 3,110 | Sunk |\n\nThe highlighted cells of the table are: [HIGHLIGHTED_BEGIN] [11,000], [Sunk], [8,096], [Damaged] [HIGHLIGHTED_END] What happened to the two heaviest ships Gerhard Bigalk attacked?
\end{lstlisting}

\textbf{Instruction:}
\begin{lstlisting}
This is a free-form table question answering task. The goal for this task is to answer the given question based on the given table and the highlighted cells.
\end{lstlisting}
\textbf{Output:}
\begin{lstlisting}
Gerhard Bigalk damaged one ship of 8,096 GRT, and sunk one warship of 11,000 tons.
\end{lstlisting}

\subsection{TabFact}
\textbf{Input:}
\begin{lstlisting}
[TLE] The table caption is about tony lema. [TAB] | tournament | wins | top - 5 | top - 10 | top - 25 | events | cuts made [SEP] | masters tournament | 0 | 1 | 2 | 4 | 4 | 4 | [SEP] | us open | 0 | 2 | 3 | 4 | 6 | 5 | [SEP] | the open championship | 1 | 2 | 2 | 2 | 3 | 3 | [SEP] | pga championship | 0 | 0 | 1 | 2 | 5 | 4 | [SEP] | totals | 1 | 5 | 8 | 12 | 18 | 16 |\n\nThe statement is:  <tony lema be in the top 5 for the master tournament , the us open , and the open championship>. Is it entailed or refuted by the table above?
\end{lstlisting}
\textbf{Instruction:}
\begin{lstlisting}
This is a table fact verification task. The goal of this task is to distinguish whether the given statement is entailed or refuted by the given table.
\end{lstlisting}
\textbf{Output:}
\begin{lstlisting}
entailed
\end{lstlisting}


\subsection{ToTTo}
\textbf{Input:}
\begin{lstlisting}
<page_title> List of Governors of South Carolina </page_title> <section_title> Governors under the Constitution of 1868 </section_title> <table> <cell> 76 <col_header> # </col_header> <col_header> 74 </col_header> <col_header> 75 </col_header> </cell> <cell> Daniel Henry Chamberlain <col_header> Governor </col_header> <row_header> 76 </row_header> </cell> <cell> December 1, 1874 <col_header> Took Office </col_header> <row_header> 76 </row_header> </cell> </table>
\end{lstlisting}
\textbf{Instruction:}
\begin{lstlisting}
This is a highlighted cells description task. The goal of this task is to generate the language description given table cells.
\end{lstlisting}
\textbf{Output:}
\begin{lstlisting}
Daniel Henry Chamberlain was the 76th Governor of South Carolina from 1874.
\end{lstlisting}


\subsection{Beer}
\textbf{Input:}
\begin{lstlisting}
Beer A is:\n|name|factory|\n|---|---|\n|Sierra Amber Ale|Silver Peak Restaurant \& Brewery|\n\nBeer B is:\n|name|factory|\n|---|---|\n|Sierra Andina Alpamayo Amber Ale|Sierra Andina|    
\# Task Description: Please determine whether Beer A and Beer B refer to the same entity or not.
\end{lstlisting}
\textbf{Instruction:}
\begin{lstlisting}
You are a helpful assistant that specializes in tables.\n Your final answer should be \'Yes\' or \'No\'. Return the final result as JSON in the format \{"answer": "<Yes or No>"\}.  Let's think step by step and show your reasoning before showing the final result.
\end{lstlisting}
\textbf{Output:}
\begin{lstlisting}
\{"answer": "No"\}
\end{lstlisting}

\subsection{TabB\textsubscript{eval}}

\textbf{Input:}
\begin{lstlisting}
Read the table below in JSON format:\n[TABLE] \n\{"columns": ["index", "organization", "year", "rank", "out of"], "data": [["bribe payers index", "transparency international", 2011, 19, 28], ["corruption perceptions index", "transparency international", 2012, 37, 176], ["democracy index", "economist intelligence unit", 2010, 36, 167], ["ease of doing business index", "world bank", 2012, 16, 185], ["economic freedom index", "fraser institute", 2010, 15, 144], ["economic freedom index", "the heritage foundation", 2013, 20, 177], ["global competitiveness report", "world economic forum", 20122013, 13, 144], ["global peace index", "institute for economics and peace", 2011, 27, 153], ["globalization index", "at kearney / foreign policy magazine", 2006, 35, 62], ["press freedom index", "reporters without borders", 2013, 47, 179], ["property rights index", "property rights alliance", 2008, 28, 115]]\}\n\nLet\'s get start!\nQuestion: What is the average rank of the indices published by Transparency International?
\end{lstlisting}
\textbf{Instruction:}
\begin{lstlisting}
You are a helpful assistant that specializes in tables.\nYou are a table analyst. Your task is to answer questions based on the table content.\n\n\nThe answer should follow the format below:\n[Answer Format]\nFinal Answer: AnswerName1, AnswerName2...\n\nEnsure the final answer format is the last output line and can only be in the "Final Answer: AnswerName1, AnswerName2..." form, no other form. Ensure the "AnswerName" is a number or entity name, as short as possible, without any explanation.\n\n\nGive the final answer to the question directly without any explanation.
\end{lstlisting}
\textbf{Output:}
\begin{lstlisting}
28
\end{lstlisting}



\subsection{MMLU}
\textbf{Input:}
\begin{lstlisting}
{5-shot examples}
Find the degree for the given field extension Q(sqrt(2), sqrt(3), sqrt(18)) over Q.	
\nA. 0\nB. 4\nC. 2\nD. 6\nAnswer:
\end{lstlisting}
\textbf{Instruction:}
\begin{lstlisting}
The following are multiple choice questions (with answers) about abstract algebra.\n\n
\end{lstlisting}
\textbf{Output:}
\begin{lstlisting}
B
\end{lstlisting}

\subsection{IFEval}
\textbf{Input:}
\begin{lstlisting}
Can you help me make an advertisement for a new product? It's a diaper that's designed to be more comfortable for babies and I want the entire output in JSON format.
\end{lstlisting}
\textbf{Instruction:}
\begin{lstlisting}
You are a helpful assistant.
\end{lstlisting}
\textbf{Output:}
\begin{lstlisting}
[JSON formatted answer]
\end{lstlisting}


\end{document}
